\documentclass[a4paper,12pt]{scrreprt}
\usepackage[utf8]{inputenc}
\usepackage[T1]{fontenc}
\usepackage[german]{babel}
\usepackage{amsmath}
\usepackage{amssymb}
\usepackage{amsthm}
\usepackage{mathtools}
\usepackage{geometry}
\usepackage[hidelinks]{hyperref}
\usepackage{listings}
\usepackage{xcolor}
\usepackage{enumitem}
\usepackage{booktabs}
\usepackage{microtype}
\usepackage{float}
\usepackage{scrhack}
\usepackage{rotating}
\usepackage[format=plain,justification=raggedright,singlelinecheck=false]{caption}

\usepackage{algorithm}
\usepackage{algpseudocode}
\usepackage{tikz}
\usepackage{pgfplots}
\usepackage{siunitx}
\usepackage{mdframed}
\usepackage{cleveref}

\usepackage{array}
\usepackage{multirow}

\pgfplotsset{compat=1.18}
\usetikzlibrary{arrows.meta,positioning,shapes.geometric,calc,fit,decorations.pathmorphing}

% Fix duplicate hyperref anchors for algorithm line numbers
\makeatletter
\renewcommand{\theHALG@line}{\thealgorithm.\arabic{ALG@line}}
\makeatother

\geometry{margin=2.5cm}

% ---------- Farben ----------
\definecolor{mainblue}{RGB}{25, 70, 140}
\definecolor{accentred}{RGB}{180, 50, 30}
\definecolor{lightgray}{RGB}{245, 245, 248}
\definecolor{darkgray}{RGB}{60, 60, 70}
\definecolor{darkgreen}{RGB}{30, 100, 50}
\definecolor{codebg}{RGB}{250, 250, 252}

\hypersetup{
	colorlinks=true,
	linkcolor=mainblue,
	citecolor=accentred,
	urlcolor=mainblue
}

% ---------- Theorem-Umgebungen ----------
\theoremstyle{definition}
\newtheorem{definition}{Definition}[chapter]
\newtheorem{satz}{Satz}[chapter]
\newtheorem{lemma}{Lemma}[chapter]
\newtheorem{korollar}{Korollar}[chapter]
\newtheorem{proposition}{Proposition}[chapter]
\newtheorem{bemerkung}{Bemerkung}[chapter]
\newtheorem{beispiel}{Beispiel}[chapter]

% ---------- Listings ----------
\lstset{
	basicstyle=\ttfamily\small,
	keywordstyle=\color{mainblue}\bfseries,
	commentstyle=\color{darkgreen}\itshape,
	stringstyle=\color{accentred},
	numbers=left,
	numberstyle=\tiny\color{darkgray},
	breaklines=true,
	frame=single,
	backgroundcolor=\color{codebg},
	rulecolor=\color{lightgray},
	literate={ö}{{\"{o}}}1{ä}{{\"{a}}}1{ü}{{\"{u}}}1
	{Ö}{{\"{O}}}1{Ä}{{\"{A}}}1{Ü}{{\"{U}}}1{ß}{{\ss}}1
}

\setlist[itemize,1]{label=--}

% ---------- Operatoren ----------
\DeclareMathOperator{\Var}{Var}
\DeclareMathOperator{\E}{\mathbb{E}}
\DeclareMathOperator{\Prob}{\mathbb{P}}
\DeclareMathOperator{\rank}{rank}
\DeclareMathOperator{\trace}{tr}
\DeclareMathOperator{\spec}{\sigma}

% ============================================================
\title{%
	\textbf{Betriebssystem-Initialisierung und\\
		Echtzeit-Scheduling mit EDF$^+$}\\[0.4em]
	\large Hardware-Grundlagen, Boot-Sequenz, PYMCA-8-Mehrkernarchitektur\\
	und formale Korrektheitsnachweise
}
\author{Stephan Epp}
\date{11. April 2026}
% ============================================================

\begin{document}
	
	\maketitle
	\tableofcontents
	\newpage
	
	% ==============================================================
	%  KAPITEL 1 – EINLEITUNG
	% ==============================================================
	\chapter{Einleitung und Motivation}
	
	Diese Arbeit präsentiert ein umfassendes Konzept für ein neues Betriebssystem, das
	auf dem \textbf{EDF$^+$}-Scheduling-Algorithmus \cite{Epp2026edfplus} als zentralem
	Paradigma aufbaut. Ausgehend von den physikalischen Grundlagen digitaler Hardware --
	R/L/C-Schaltkreisen, Analog-Digital-Wandlern und Signaltheorie -- über die
	vollständige x86-64-Boot-Sequenz (Reset-Vektor, BIOS/UEFI-POST, ACPI,
	Bootloader-Kernel-Ãœbergang) bis hin zur Prozessor-Initialisierung in
	Real-Modus, Protected-Modus und Long-Modus werden alle wesentlichen Schichten
	eines modernen Betriebssystems formal behandelt.
	
	Der \textbf{EDF$^+$}-Scheduler erweitert klassisches Earliest-Deadline-First
	durch einen dynamischen Penalty-Mechanismus für Blocking-Ereignisse und
	wird als optimaler Algorithmus für die Hardware-Initialisierungssequenz
	eingesetzt. Darauf aufbauend beschreibt die Arbeit die
	\textbf{PYMCA-8}-Mehrkernarchitektur \cite{Epp2026pymca8} mit acht Kernen,
	Python-Speichermodell, IoFET-basiertem Dynamic-Voltage-and-Frequency-Scaling
	(DVFS) sowie das stochastische Archiv-Problem \cite{Epp2026archv}.
	
	Weitere Kapitel behandeln graphentheoretische OS-Verifikation mittels
	Subgraph Algorithmus \cite{Epp2026subgraph} und Boolescher Matrixmultiplikation
	\cite{Epp2026boolmm}, Lyapunov-Stabilität von Systemzustandsmatrizen
	\cite{Epp2026sysstate}, Dateisystem-Deduplizierung \cite{Epp2026dusgraph},
	AUTOSAR-Treiber\-architektur \cite{Epp2026autsre}, EtherCAT \cite{Epp2026ethercate},
	Energiemanagement mit IoFET-Sigmoid-DVFS und Li-Ionen-Akku-Modell
	\cite{Epp2026liionp,Epp2026iono} sowie Post-Quanten-Kryptographie
	\cite{Epp2026pqc}.
	
	Im Kapitel Neue Erkenntnisse werden sechs originäre Sätze formal bewiesen:
	die OS-Initialisierungssequenz als EDF$^+$-Task-System (Satz~14.1),
	Subgraph-basierte Treiber-Kompatibilität in $O(n^3)$ (Satz~14.2),
	stochastische Bootzeit-Schranke mit Hoeffding-Ungleichung (Satz~14.3),
	Lyapunov-Stabilität des Betriebssystems (Satz~14.4),
	Optimalität der IoFET-DVFS-Kurve beim Booten (Korollar~14.5) sowie
	Post-Quanten-sichere Speicheradressierung via LWE (Satz~14.6).
	
	\section{Warum ein neues Betriebssystem?}
	
	Moderne Betriebssysteme wie Linux, Windows oder macOS sind über Jahrzehnte
	gewachsen und tragen erheblichen historischen Ballast. Ihre Scheduling-Strategien
	stammen aus einer Ära, in der Blocking-Ereignisse eher die Ausnahme als die Regel
	waren und Echtzeit-Anforderungen allenfalls in Spezialsystemen auftraten. Heute
	hingegen sind eingebettete Systeme in Fahrzeugen (AUTOSAR, ISO~26262 \cite{ISO26262}),
	Industrierobotern (EtherCAT \cite{IEC61158}) und KI-Beschleunigern allgegenwärtig --
	allesamt Domänen, die harte Echtzeit-Garantien, vorhersagbaren Energieverbrauch
	und formale Korrektheitsnachweise verlangen.
	
	Das in dieser Arbeit beschriebene Betriebssystem \textbf{BS/EDF$^+$} setzt
	EDF$^+$-Scheduling \cite{Epp2026edfplus} als einheitliches Paradigma über alle
	Systemschichten ein: von der ersten Instruktion nach dem Reset bis zum Scheduling
	der Benutzeranwendungen. Dadurch entsteht eine vertikal integrierte,
	formal beweisbare Systemarchitektur, die bisherige Ansätze in drei wesentlichen
	Punkten übertrifft:
	
	\begin{enumerate}
		\item \textbf{Formale Korrektheit:} Jede Hardware-Initialisierungsphase ist
		ein Task im Sinne von EDF$^+$; die Schedulierbarkeit ist beweisbar.
		\item \textbf{Vorhersagbares Energieverhalten:} IoFET-Sigmoid-DVFS liefert
		die energieoptimale Spannung-Frequenz-Kurve beim Booten (Korollar~14.5).
		\item \textbf{Quantensicherheit:} Die physikalische Speicheradressierung
		ist durch LWE-basierte Gitterkryptographie quantensicher (Satz~14.6).
	\end{enumerate}
	
	\section{Benutzerperspektive}
	
	Der Benutzer erlebt das neue Betriebssystem als eine Folge sichtbarer
	Zustände, beginnend beim Einschalten:
	
	\begin{enumerate}
		\item \textbf{POST-Meldungen} (ca.\ 0--500\,ms): Speicher-Check, CPU-Identifikation.
		\item \textbf{Bootlogo} (ca.\ 500\,ms--2\,s): Grafiktreiber-Init, DVFS-Aktivierung.
		\item \textbf{Anmeldebildschirm} (ca.\ 2--4\,s): Dateisystem gemountet, Netzwerk bereit.
		\item \textbf{Desktop} (ca.\ 4--6\,s): Alle Treiber aktiv, EDF$^+$-Scheduler läuft.
	\end{enumerate}
	
	Satz~14.3 (Stochastische Bootzeit-Schranke) garantiert formal, dass die
	Gesamtbootzeit mit hoher Wahrscheinlichkeit unter einer berechneten Grenze
	$t_{\max}$ liegt.
	
	\section{Wissenschaftliche Beiträge}
	
	Die Hauptbeiträge dieser Arbeit sind:
	
	\begin{itemize}
		\item \textbf{Theoretisch:} Sechs neue Sätze mit vollständigen formalen Beweisen
		(Kapitel~14), die EDF$^+$, Subgraph Algorithmus, stochastische Analyse,
		Lyapunov-Stabilität, DVFS-Optimalität und Post-Quanten-Sicherheit
		zu einem kohärenten Framework verbinden.
		\item \textbf{Architektonisch:} Vollständige Beschreibung aller OS-Schichten
		(Kapitel~2--13) auf Basis der referenzierten eigenen Arbeiten.
		\item \textbf{Algorithmisch:} Integration des Subgraph Algorithmus
		\cite{Epp2026subgraph} in die Treiber-Kompatibi\-litätsprüfung (O(n$^3$)).
	\end{itemize}
	
	\section{Aufbau der Arbeit}
	
	Kapitel~2 beschreibt die Computerhardware-Grundlagen; Kapitel~3 die
	physikalischen Grundlagen (R/L/C, Fourier, ADC/DAC). Kapitel~4 behandelt die
	Boot-Sequenz, Kapitel~5 den EDF$^+$-Scheduler im Detail,
	Kapitel~6 die PYMCA-8-Architektur. Kapitel~7 analysiert das Archiv- und
	Speicherverwaltungsproblem. Kapitel~8--13 decken Graphenverifikation,
	Systemzustände, Dateisystem, Treiber, Energie und Kryptographie ab.
	Kapitel~14 enthält die neuen Beweise, Kapitel~15 Zusammenfassung und Ausblick.
	
	% ==============================================================
	%  KAPITEL 2 – COMPUTERHARDWARE-GRUNDLAGEN
	% ==============================================================
	\chapter{Computerhardware-Grundlagen}
	
	\section{Ãœberblick: Hardwarekomponenten}
	
	Ein moderner Computer besteht aus einer hierarchisch organisierten Menge von
	Hardwarekomponenten, die über standardisierte Bussysteme kommunizieren.
	Abbildung~\ref{fig:hw-overview} zeigt schematisch die wesentlichen Elemente und
	deren Verbindungsstruktur über den PCIe-Root-Complex und den Speicher-Controller
	als zentralen Verbindungsknoten.

	\begin{sidewaysfigure}
		\centering
		\begin{tikzpicture}[
			node distance=1.6cm and 2.4cm,
			block/.style={rectangle, draw=mainblue, fill=mainblue!12, rounded corners=4pt,
				text width=2.6cm, align=center, minimum height=0.85cm, font=\small},
			sblock/.style={rectangle, draw=accentred, fill=accentred!10, rounded corners=4pt,
				text width=2.4cm, align=center, minimum height=0.75cm, font=\small},
			bus/.style={-{Stealth[length=5pt]}, thick, mainblue},
			dbus/.style={latex-latex, thick, darkgray, dashed}
		]
			\node[block] (cpu) {\textbf{CPU}\\8 Kerne\\x86-64};
			\node[block, right=2.5cm of cpu] (mc) {\textbf{Speicher-\\Controller}};
			\node[block, right=2.5cm of mc] (ram) {\textbf{DDR5-DRAM}\\16--512\,GB};

			\node[block, below=1.8cm of cpu] (pcie) {\textbf{PCIe Root\\Complex}\\Gen\,5};
			\node[sblock, below left=1.5cm and 0.2cm of pcie] (gpu) {\textbf{GPU}\\PCIe$\times$16};
			\node[sblock, below=1.5cm of pcie] (nvme) {\textbf{NVMe SSD}\\PCIe$\times$4};
			\node[sblock, below right=1.5cm and 0.2cm of pcie] (nic) {\textbf{NIC}\\100\,GbE};

			\node[sblock, right=2.4cm of pcie] (usb) {\textbf{USB\,3.2}\\Hub};
			\node[sblock, left=2.4cm of pcie] (sata){\textbf{SATA\,III}\\Controller};

			\draw[bus] (cpu) -- node[above,font=\scriptsize]{QPI/UPI} (mc);
			\draw[bus] (mc) -- node[above,font=\scriptsize]{64\,bit\,bus} (ram);
			\draw[dbus] (cpu) -- (pcie);
			\draw[bus] (pcie) -- (gpu);
			\draw[bus] (pcie) -- (nvme);
			\draw[bus] (pcie) -- (nic);
			\draw[bus] (pcie) -- (usb);
			\draw[bus] (pcie) -- (sata);
		\end{tikzpicture}
		\caption{Schematische Darstellung der Systemhardware. Der Speicher-Controller
			verbindet die CPU direkt mit dem DDR5-DRAM (On-Die-Integration bei modernen
			Prozessoren). Der PCIe-Root-Complex verteilt die Bandbreite an alle
			Peripheriegeräte; GPU und NVMe-SSD erhalten die breitesten Lanes.}
		\label{fig:hw-overview}
	\end{sidewaysfigure}

	Die Kommunikation zwischen CPU und DRAM erfolgt über einen
	\emph{integrierten Speicher-Controller} (IMC), der bei modernen Intel- und
	AMD-Prozessoren direkt im Die sitzt. Dadurch sinkt die Zugriffslatenz gegenüber
	externen Speicher-Controllern (Northbridge) von über 100\,ns auf etwa 65\,ns.
	Der PCIe-Root-Complex ist logisch ein Switch, dessen Bandbreite von GPU und
	NVMe-SSDs gemeinsam genutzt wird; BS/EDF$^+$ modelliert diese Bandbreite
	als geteilte Ressource innerhalb des EDF$^+$-Task-Systems (Kapitel~5).

	\begin{definition}[Computersystem]
		Ein \textbf{Computersystem} ist ein Tupel
		\[
		\mathcal{H} = (P, M, S, I, B)
		\]
		mit Prozessormenge~$P$, Speicherhierarchie~$M$, Speichergeräten~$S$,
		Ein-/Ausgabe-Geräten~$I$ und Busnetzwerk~$B$.
	\end{definition}
	
	\section{Prozessor (CPU)}
	
	\subsection{Register und Ausführungsmodell}
	
	Der x86-64-Prozessor besitzt 16 Allzweckregister (RAX, RBX, \ldots{}, R15),
	sechs Segmentregister (CS, DS, ES, FS, GS, SS), den Instruction Pointer (RIP),
	das Flags-Register (RFLAGS) sowie eine Vielzahl von Steuerregistern (CR0--CR8),
	Model-Specific-Registern (MSR) und den APIC-Basis-Registern \cite{Intel2023}.
	
	\begin{definition}[Prozessormodus]
		Ein x86-64-Prozessor operiert in einem von vier Hauptmodi:
		\begin{enumerate}
			\item \textbf{Real-Modus} (16-Bit, Adressraum $2^{20}$\,Byte = 1\,MB)
			\item \textbf{Protected-Modus} (32-Bit, Adressraum $2^{32}$\,Byte = 4\,GB)
			\item \textbf{Long-Modus} (64-Bit, Adressraum $2^{48}$\,Byte = 256\,TB effektiv)
			\item \textbf{System-Management-Modus} (SMM, gesonderter Adressraum SMRAM)
		\end{enumerate}
	\end{definition}
	
	\subsection{Instruktions-Pipeline}

	Moderne x86-64-Prozessoren nutzen eine \emph{Out-of-Order-Pipeline} mit bis zu
	20 Stufen. Das Grundprinzip lässt sich anhand einer vereinfachten
	5-Stufen-In-Order-Pipeline verdeutlichen, wie sie in eingebetteten Prozessoren
	(z.\,B.\ ARM Cortex-M) vorkommt und als konzeptuelles Modell für
	Latenz-Analysen dient.

	\begin{figure}[H]
		\centering
		\begin{tikzpicture}[
			stage/.style={rectangle, draw=mainblue, fill=mainblue!15, minimum width=2cm,
				minimum height=0.9cm, align=center, font=\small\bfseries, rounded corners=3pt},
			arr/.style={{Stealth[length=6pt]}-{Stealth[length=6pt]}, thick, darkgray}
		]
			\foreach \name/\label/\i in {
				IF/Fetch/0, ID/Decode/1, EX/Execute/2, MEM/Memory/3, WB/Writeback/4}
			{
				\node[stage] at (\i*2.4,0) (\name) {\name\\\scriptsize\label};
			}
			\draw[-{Stealth}, thick, mainblue] (IF) -- (ID);
			\draw[-{Stealth}, thick, mainblue] (ID) -- (EX);
			\draw[-{Stealth}, thick, mainblue] (EX) -- (MEM);
			\draw[-{Stealth}, thick, mainblue] (MEM) -- (WB);

			% Hazard annotation
			\draw[-{Stealth}, accentred, thick, bend left=45]
				(ID.north) to node[above, font=\scriptsize, accentred]{Data-Hazard} (EX.north);
			\draw[-{Stealth}, darkgreen, thick, bend right=-40]
				(EX.south) to node[below, font=\scriptsize, darkgreen]{Forwarding} (ID.south);
		\end{tikzpicture}
		\caption{Vereinfachte 5-Stufen-Pipeline (IF--ID--EX--MEM--WB). Ein
			\emph{Data-Hazard} entsteht, wenn eine Stufe auf ein Ergebnis wartet,
			das noch in einer vorangehenden Stufe berechnet wird. \emph{Forwarding}
			(gestrichelte grüne Linie) leitet das Ergebnis direkt aus der EX-Stufe
			zurück, ohne einen Stall einzufügen. BS/EDF$^+$ berücksichtigt
			Pipeline-Stall-Wahrscheinlichkeiten bei der Worst-Case-Execution-Time
			(WCET)-Schätzung der EDF$^+$-Tasks.}
		\label{fig:pipeline}
	\end{figure}

	Die \emph{Out-of-Order}-Ausführung moderner Prozessoren (Intel Raptor Lake,
	AMD Zen~4) nutzt einen \emph{Reorder Buffer} (ROB) mit bis zu 512 Einträgen.
	Instruktionen werden nach ihrer Fertigstellung in Programm\-reihenfolge
	\emph{retired}, sodass der architektursichtbare Zustand stets konsistent bleibt.
	Für Echtzeit-Scheduling nach EDF$^+$ ist maßgeblich, dass die maximale
	Instruktionslatenz durch den \emph{Worst-Case Execution Path} (WCEP)
	beschreibbar ist -- ein Pfad, auf dem alle Branches die ungünstigste
	Richtung einschlagen und alle Cache-Zugriffe misslingen.

	\begin{definition}[WCET]
		Die \textbf{Worst-Case Execution Time} (WCET) eines Tasks $T_i$ ist
		\[
		C_i^{\max} = \max_{\text{Eingabe}} \bigl\{\text{Cycles}(T_i, \text{Eingabe})\bigr\}
		\cdot T_{\text{Clk}},
		\]
		wobei $T_{\text{Clk}} = 1/f_{\text{CPU}}$ die Zykluszeit ist. EDF$^+$
		verwendet $C_i = C_i^{\max}$ zur sicheren Schedulierbarkeitsanalyse.
	\end{definition}

	\subsection{Cache-Kohärenz und MESI-Protokoll}

	Auf Mehr\-kern-Systemen wie PYMCA-8 müssen alle Kerne ein konsistentes Bild
	des Hauptspeichers haben. Das \textbf{MESI-Protokoll} definiert vier
	Zustände für jede Cache-Line (64\,Byte):

	\begin{center}
		\begin{tabular}{lll}
			\toprule
			Zustand & Bedeutung & Bedingung \\
			\midrule
			\textbf{M} (Modified)  & Geändert, nur in diesem Cache & Dirty, exklusiv \\
			\textbf{E} (Exclusive) & Sauber, nur in diesem Cache    & Clean, exklusiv \\
			\textbf{S} (Shared)    & Sauber, ggf.\ in mehreren Caches & Clean, geteilt \\
			\textbf{I} (Invalid)   & Ungültig                        & Muss neu geladen werden \\
			\bottomrule
		\end{tabular}
	\end{center}

	Bei einem \emph{Write} auf eine Shared-Line senden alle anderen Kerne ein
	\emph{Invalidate}-Signal, bevor der schreibende Kern seine Kopie auf Modified
	setzt. Dieser \emph{Cache-Invalidierungsoverhead} ist für EDF$^+$ relevant:
	Bei Task-Migrationen zwischen PYMCA-8-Kernen kann das Neueinladen von
	Cache-Lines zusätzliche Zyklen kosten, die in die WCET-Schätzung einfließen.

	\subsection{Cache-Hierarchie}

	\begin{definition}[Speicherhierarchie]
		Die \textbf{Speicherhierarchie} $M = (M_1, M_2, \ldots, M_k)$ ist geordnet
		nach Zugriffszeit $t_i$ und Kapazität $C_i$ mit $t_1 < t_2 < \cdots < t_k$
		und $C_1 < C_2 < \cdots < C_k$.
	\end{definition}
	
	Typische Werte für aktuelle x86-64-Prozessoren:
	
	\begin{center}
		\begin{tabular}{llll}
			\toprule
			Ebene & Typ & Kapazität & Latenz \\
			\midrule
			L1-D/I & SRAM & 32--64\,KB & 4 Zyklen \\
			L2     & SRAM & 256\,KB--1\,MB & 12 Zyklen \\
			L3     & SRAM & 4--64\,MB & 40 Zyklen \\
			DRAM   & DDR5 & 8--512\,GB & 80--100\,ns \\
			NVMe   & Flash & 500\,GB--8\,TB & 20--100\,$\mu$s \\
			\bottomrule
		\end{tabular}
	\end{center}
	
	\subsection{CPUID und MSR}
	
	Die \texttt{CPUID}-Instruktion liefert Prozessor-Identifikation und
	Feature-Flags. Wichtige Blätter:
	\begin{itemize}
		\item \texttt{EAX=0x01}: Family/Model/Stepping, Standard-Features (SSE, AVX, AES-NI)
		\item \texttt{EAX=0x07}: Extended Features (AVX-512, SHA)
		\item \texttt{EAX=0x80000001}: Extended Features AMD/Intel (Long-Modus-Bit)
	\end{itemize}
	
	Model-Specific-Register werden via \texttt{RDMSR}/\texttt{WRMSR}
	gelesen/geschrieben. Für EDF$^+$-Sche\-duling relevant sind:
	MSR \texttt{0x0C0000100} (FS.base), \texttt{0x0C0000101} (GS.base) für
	thread-lokale Speicherung sowie \texttt{IA32\_PERF\_CTL} (0x199) für DVFS.
	
	\section{Speicher (RAM)}
	
	Moderner DDR5-DRAM ist organisiert in Channels, DIMMs, Ranks, Banks und
	Rows/Columns. Für das Betriebssystem sind folgende physikalische Adressen
	nach dem POST reserviert:
	
	\begin{center}
		\begin{tabular}{ll}
			\toprule
			Adressbereich & Verwendung \\
			\midrule
			\texttt{0x00000000}--\texttt{0x0009FFFF} & Konventioneller Speicher (640\,KB) \\
			\texttt{0x000A0000}--\texttt{0x000FFFFF} & ROM/VGA-BIOS, Legacy \\
			\texttt{0x00100000}--\texttt{Ende DRAM} & Extended Memory (OS-nutzbar) \\
			\bottomrule
		\end{tabular}
	\end{center}
	
	\section{GPU und Grafiksubsystem}
	
	Moderne Grafikprozessoren (NVIDIA H100, AMD RX 7900) kommunizieren über
	PCIe Gen~4/5 mit dem Host. Das BS/EDF$^+$ initialisiert die GPU im
	Bootprozess als EDF$^+$-Task $T_{\text{GPU}}$ mit Deadline $D_{\text{GPU}}$,
	die durch das Booten des grafischen Terminals bedingt ist.
	Für HPC-Workloads unterstützt PYMCA-8 \cite{Epp2026pymca8} CUDA/ROCm-Integration
	mit IoFET-DVFS auch auf der GPU-Seite.
	
	\section{PCIe, NVMe und SATA}
	
	PCI Express (PCIe Gen~5: 32\,GT/s pro Lane, $\times$16: 64\,GB/s bidirektional)
	verbindet GPU, NVMe-SSDs und Netzwerkkarten mit dem Prozessor über den
	Root Complex. Der BS/EDF$^+$-Boot-Scheduler behandelt jede PCIe-Geräteinitiali\-sierung
	als separaten Task mit angemessener Deadline.
	
	NVMe-SSDs verwenden bis zu 64\,000 I/O-Queues (im Gegensatz zu AHCI: 1 Queue,
	32 Commands), wodurch parallele Initialisierung im EDF$^+$-Framework möglich ist.
	
	\section{USB und Netzwerk}
	
	USB~3.2 Gen~2×2 erreicht 20\,Gbit/s. Netzwerkkarten (10-GbE, 100-GbE) werden
	über UEFI PXE-Boot oder nach dem Kernel-Start über DHCP initialisiert.
	Für eingebettete Systeme ist EtherCAT \cite{IEC61158} relevant, das in
	Kapitel~11 detailliert behandelt wird.
	
	% ==============================================================
	%  KAPITEL 3 – PHYSIKALISCHE GRUNDLAGEN
	% ==============================================================
	\chapter{Physikalische Grundlagen}
	
	\section{R/L/C-Schaltkreise und Zeitkonstanten}

	Digitale Hardware basiert auf Schaltkreisen aus Widerständen ($R$),
	Induktivitäten ($L$) und Kapazitäten ($C$). Die Zeitkonstante bestimmt
	Signalverzögerungen und damit Taktlimits \cite{Epp2026physd}.

	\begin{definition}[RC-Zeitkonstante]
		Für einen einfachen RC-Tiefpass gilt die Zeitkonstante
		\[
		\tau_{RC} = R \cdot C.
		\]
		Die Ausgangsspannung bei Sprungeingang $U_0$ ist
		\[
		u_C(t) = U_0 \left(1 - e^{-t/\tau_{RC}}\right).
		\]
	\end{definition}

	Abbildung~\ref{fig:rc-charge} zeigt die normierte Ladekurve
	$u_C(t)/U_0$ für verschiedene Zeitkonstanten. Für Prozessortakte
	der heutigen Generation müssen Bypass-Kondensatoren mit sehr kleiner
	Impedanz ($Z < 1\,\Omega$ bei der CPU-Taktfrequenz) eingesetzt werden,
	da bei $f_{\text{CPU}} = 5\,\text{GHz}$ und $\tau_{RC} = 200\,\text{ps}$
	der Spannungsabfall bei transienten Laststößen bereits die
	Stabilitätsschwelle des Prozessors erreichen kann.

	\begin{figure}[H]
		\centering
		\begin{tikzpicture}
			\begin{axis}[
				width=11cm, height=6cm,
				xlabel={Zeit $t/\tau_{RC}$},
				ylabel={$u_C/U_0$},
				xmin=0, xmax=5, ymin=0, ymax=1.05,
				grid=major, grid style={lightgray!60},
				xtick={0,1,2,3,4,5},
				ytick={0,0.2,0.4,0.6,0.8,1.0},
				legend pos=south east,
				cycle list name=color list,
			]
				\addplot[mainblue, thick, domain=0:5, samples=200]
					{1 - exp(-x)};
				\addlegendentry{Ladung $u_C(t) = U_0(1-e^{-t/\tau})$}

				\addplot[accentred, thick, dashed, domain=0:5, samples=200]
					{exp(-x)};
				\addlegendentry{Entladung $u_C(t) = U_0 e^{-t/\tau}$}

				\addplot[darkgray, thin, dashed] coordinates {(1,0)(1,0.632)};
				\addplot[darkgray, thin, dashed] coordinates {(0,0.632)(1,0.632)};
				\node[font=\scriptsize] at (axis cs:1.3,0.55) {$(1,\, 0.632)$};
			\end{axis}
		\end{tikzpicture}
		\caption{Normierte Lade- und Entladekurve eines RC-Glieds. Nach einer
			Zeitkonstante $\tau_{RC}$ hat der Kondensator $63.2\,\%$ seiner
			Endspannung erreicht (roter Punkt). Im Kontext von Taktleitungen auf
			Leiterplatten entspricht die Zeitkonstante der Flankenzeit des
			Taktsignals; für $f_{\text{CPU}} = 5\,\text{GHz}$ muss
			$\tau_{RC} \ll 100\,\text{ps}$ gelten.}
		\label{fig:rc-charge}
	\end{figure}

	\begin{definition}[RLC-Schwingkreis]
		Ein serieller RLC-Schwingkreis hat die charakteristische Gleichung
		\[
		L\ddot{i} + R\dot{i} + \frac{1}{C}i = \frac{d u_e}{dt}.
		\]
		Die Eigenfrequenz beträgt $\omega_0 = 1/\sqrt{LC}$ und der
		Dämpfungsgrad $\delta = R/(2L)$. Das System ist:
		\begin{itemize}
			\item \emph{überdämpft} falls $\delta > \omega_0$,
			\item \emph{kritisch gedämpft} falls $\delta = \omega_0$,
			\item \emph{schwingend} falls $\delta < \omega_0$.
		\end{itemize}
	\end{definition}
	
	Für Prozessortakte mit $f_{\text{CPU}} = 3$--$5\,\text{GHz}$ müssen
	Leitungsinduktivitäten im Bereich $L < 1\,\text{nH}$ und
	Bypass-Kondensatoren $C \approx 100\,\mu\text{F}$ dimensioniert werden,
	um $\tau_{LC} = \sqrt{LC} < 10\,\text{ps}$ zu gewährleisten.
	
	\section{Fourier-Analyse und Signaldarstellung}
	
	\begin{definition}[Fourier-Transformation]
		Die \textbf{kontinuierliche Fourier-Transformation} eines Signals $x(t)$ ist
		\[
		X(f) = \int_{-\infty}^{\infty} x(t)\, e^{-j2\pi f t}\, dt.
		\]
		Die Rücktransformation lautet:
		\[
		x(t) = \int_{-\infty}^{\infty} X(f)\, e^{j2\pi f t}\, df.
		\]
	\end{definition}
	
	Für Rechteckpulse der Breite $\tau$ und Periode $T$ (Taktsignal) gilt:
	\[
	X_n = \frac{\tau}{T} \cdot \text{sinc}\!\left(\frac{n\tau}{T}\right),
	\]
	wobei $\text{sinc}(x) = \sin(\pi x)/(\pi x)$. Die Bandbreite eines
	Takt\-signals der Frequenz $f_0$ beträgt typischerweise $B = 5f_0$,
	was für 5\,GHz also 25\,GHz Bandbreite erfordert -- relevant für
	PCB-Design und Signalintegrität \cite{Epp2026signalth}.
	
	\section{Abtasttheorem und Analog-Digital-Wandlung}
	
	\begin{satz}[Nyquist-Shannon-Abtasttheorem]
		Ein bandbegrenztes Signal $x(t)$ mit $X(f) = 0$ für $|f| > f_{\max}$
		kann aus seinen Abtastwerten $x(nT_s)$ eindeutig rekonstruiert werden,
		falls die Abtastrate
		\[
		f_s = \frac{1}{T_s} \geq 2 f_{\max}.
		\]
	\end{satz}
	
	\begin{proof}
		Die Abtastung im Zeitbereich entspricht einer Faltung mit einem Dirac-Kamm
		$s(t) = \sum_n \delta(t - nT_s)$, was im Frequenzbereich eine Periodisierung
		mit Periode $f_s$ ergibt: $X_s(f) = f_s \sum_k X(f - kf_s)$. Falls
		$f_s \geq 2f_{\max}$, überlappen die periodisierten Spektren nicht, und
		$X(f)$ kann durch einen idealen Tiefpass mit Grenzfrequenz $f_s/2$
		vollständig zurückgewonnen werden.
	\end{proof}
	
	\begin{definition}[ADC-Quantisierungsrauschen]
		Ein $n$-Bit-ADC mit Eingangsbereich $[-U_{\max}, U_{\max}]$ hat die
		Quantisierungsstufe $\Delta = 2U_{\max}/2^n$ und den Signal-zu-Rausch-Abstand
		\[
		\mathrm{SNR}_{\mathrm{dB}} = 6{,}02 \cdot n + 1{,}76\,\mathrm{dB}.
		\]
	\end{definition}
	
	Für Sensorinterfaces in eingebetteten Systemen (z.\,B.\ Temperatur-, Druck-
	oder Stromsensoren in AUTOSAR-ECUs) sind typischerweise 12--16-Bit-ADCs
	mit $f_s = 1$--$100\,\mathrm{kHz}$ im Einsatz.
	
	\section{Energiebilanz digitaler Schaltkreise}
	
	Die dynamische Verlustleistung eines CMOS-Gatters beträgt:
	\[
	P_{\text{dyn}} = \alpha \cdot C_L \cdot V_{DD}^2 \cdot f,
	\]
	wobei $\alpha$ die Schaltaktivität, $C_L$ die Lastkapazität, $V_{DD}$
	die Versorgungsspannung und $f$ die Taktfrequenz ist. Dies ist die
	physikalische Grundlage des DVFS-Konzepts (Kapitel~6 und~12): durch
	simultane Reduktion von $V_{DD}$ und $f$ sinkt die Leistung kubisch.

	\section{Signalintegrität und Übertragungsleitungen}

	Bei Signalfrequenzen oberhalb von etwa $100\,\text{MHz}$ und Leitungslängen
	im Bereich der Wellenlänge muss eine Leitung als \emph{Übertragungsleitung}
	modelliert werden (Telegrafen-Gleichungen):
	\[
	\frac{\partial u}{\partial z} = -L'\frac{\partial i}{\partial t} - R' i,
	\qquad
	\frac{\partial i}{\partial z} = -C'\frac{\partial u}{\partial t} - G' u,
	\]
	mit Belagsgrößen $R'$, $L'$, $C'$, $G'$ (pro Längeneinheit).
	Der \textbf{Wellenwiderstand} beträgt
	\[
	Z_0 = \sqrt{\frac{L'}{C'}} \approx 50\,\Omega
	\]
	für typische PCB-Mikro\-streifenleitungen. Eine Fehlanpassung
	$Z_{\mathrm{Last}} \neq Z_0$ erzeugt Reflexionen, die das Setup-und-Hold-Fenster
	von Flip-Flops gefährden können.

	\begin{bemerkung}[DVFS und Signalintegrität]
		Eine DVFS-Spannungsabsenkung $V_{DD} \to V_{DD}'$ verringert die
		Spannungspegelabstände und damit das Rauschmargen-Budget (\emph{Noise Margin}).
		BS/EDF$^+$ beschränkt daher DVFS-Absenkungen auf maximal $\Delta V = 200\,\text{mV}$
		pro Takt-Domäne, um $V_{\mathrm{NM}} > 150\,\text{mV}$ zu garantieren.
	\end{bemerkung}
	
	% ==============================================================
	%  KAPITEL 4 – BOOT-SEQUENZ
	% ==============================================================
	\chapter{Boot-Sequenz}
	
	\section{Ãœberblick}

	Die Boot-Sequenz umfasst alle Schritte vom Anlegen der Versorgungsspannung
	bis zum vollständig initialisierten Kernel. Sie gliedert sich in sechs Phasen:

	\begin{enumerate}
		\item \textbf{Power-On-Reset}: Versorgungsspannung stabilisiert sich,
		Reset-Vektor wird geladen.
		\item \textbf{POST} (Power-On Self-Test): BIOS/UEFI prüft und initialisiert
		Hardware.
		\item \textbf{ACPI-Tabellen}: System-Topologie wird beschrieben.
		\item \textbf{Bootloader}: MBR oder EFI-Partition, Kernel-Image wird geladen.
		\item \textbf{Kernel-Initialisierung}: Decompress, Speicher-Setup,
		Treiber-Probe.
		\item \textbf{Init-System}: Systemd/OpenRC, Dienste werden nach EDF$^+$
		gestartet.
	\end{enumerate}

	Abbildung~\ref{fig:boot-flow} illustriert den gesamten Ablauf als
	sequentielle Zustandsmaschine. Jeder Ãœbergang zwischen Phasen ist ein
	\emph{Task-Release}-Ereignis im EDF$^+$-Framework.Der Ablauf des Algorithmus lässt sich wie folgt beschreiben: 
	Beim Anlegen der Betriebsspannung beginnt der Prozessor mit dem Ausführen von
	Mikrocode-Routinen, die den internen Zustand des Prozessors initialisieren
	(Register auf Null setzen, Fehlerflags löschen). Danach wird der erste
	Maschinenbefehl am Reset-Vektor \texttt{0xFFFFFFF0} abgerufen -- ein
	unbedingter Sprung in den BIOS/UEFI-ROM-Code. Die UEFI-Firmware durchläuft
	dann die SEC-, PEI-, DXE-, und BDS-Phasen (Abschnitt~\ref{sec:uefi_phases})
	und übergibt die Kontrolle schließlich an den Bootloader. Dieser
	dekomprimiert das Kernel-Image und springt in den Kernel-Einstiegspunkt,
	wo BS/EDF$^+$ den ersten EDF$^+$-Scheduling-Kontext aufbaut.

	\begin{figure}[H]
		\centering
		\begin{tikzpicture}[
			every node/.style={font=\small},
			phase/.style={rectangle, draw=mainblue, fill=mainblue!12, rounded corners=5pt,
				minimum width=4.4cm, minimum height=0.75cm, align=center},
			check/.style={diamond, draw=accentred, fill=accentred!10,
				minimum width=3.2cm, minimum height=0.7cm, align=center, aspect=2.2},
			arr/.style={-{Stealth[length=5pt]}, thick, darkgray},
			node distance=0.6cm
		]
			\node[phase] (por)   {Power-On-Reset\\Reset-Vektor \texttt{0xFFFFFFF0}};
			\node[phase, below=of por] (post) {BIOS/UEFI POST\\CPU, RAM, PCIe};
			\node[check, below=of post] (postok) {POST OK?};
			\node[phase, below=of postok] (acpi) {ACPI-Tabellen lesen\\MADT, SRAT, HPET};
			\node[phase, below=of acpi]  (boot)  {Bootloader (GRUB2)\\Stage 1 \textrightarrow{} 2};
			\node[phase, below=of boot]  (kdec)  {Kernel entpacken\\zstd/LZ4};
			\node[phase, below=of kdec]  (kinit) {Kernel-Init\\head\_64.S, start\_kernel()};
			\node[phase, below=of kinit] (sched) {EDF$^+$-Scheduler starten\\sched\_init()};
			\node[phase, below=of sched] (init)  {Init-System\\Dienste (EDF$^+$-Tasks)};

			\node[phase, right=2.2cm of postok, fill=accentred!15, draw=accentred]
				(err) {Fehler-LED/Beep\\Halt};

			\draw[arr] (por)   -- (post);
			\draw[arr] (post)  -- (postok);
			\draw[arr] (postok) -- node[right,font=\scriptsize]{Ja} (acpi);
			\draw[arr] (postok) -- node[above,font=\scriptsize]{Nein} (err);
			\draw[arr] (acpi)  -- (boot);
			\draw[arr] (boot)  -- (kdec);
			\draw[arr] (kdec)  -- (kinit);
			\draw[arr] (kinit) -- (sched);
			\draw[arr] (sched) -- (init);
		\end{tikzpicture}
		\caption{Vollständiges Boot-Ablaufdiagramm von BS/EDF$^+$. Jede Phase
			entspricht einem oder mehreren EDF$^+$-Tasks. Der POST-Ausgangstest
			(\emph{POST OK?}) trennt den normalen Startpfad (grüner Weg nach unten)
			vom Fehlerpfad (rechts). Ab der Kernel-Initialisierungsphase übernimmt
			der EDF$^+$-Scheduler die Kontrolle über alle weiteren Tasks.}
		\label{fig:boot-flow}
	\end{figure}

	\section{Reset-Vektor und erste Instruktion}
	
	\begin{definition}[Reset-Vektor]
		Nach einem Power-On-Reset oder Hard-Reset setzt der x86-64-Prozessor den
		Instruction Pointer auf den \textbf{Reset-Vektor}:
		\[
		\texttt{CS:IP} = \texttt{0xF000:0xFFF0} \equiv \texttt{0xFFFFFFF0}.
		\]
		An dieser Adresse befindet sich ein JMP-Befehl zum BIOS-Einstiegspunkt.
	\end{definition}
	
	Die ersten Instruktionen laufen im Real-Modus (16-Bit). Der Prozessor befindet
	sich in einem definierten Zustand: alle Register auf 0, EFLAGS = 0x00000002,
	CR0 = 0 (Protected Mode Disabled), EFLAG.IF = 0 (Interrupts disabled).
	
	\section{BIOS und UEFI POST}
	
	\subsection{Legacy-BIOS}
	
	Das legacy BIOS führt folgende Schritte aus:
	\begin{enumerate}
		\item \textbf{CPU-Test}: CPUID-Aufruf, Familien-Identifikation.
		\item \textbf{Speichertest}: Walking-1-Muster über alle DRAM-Adressen.
		\item \textbf{PCI-Enumeration}: Bus-Scan, BAR-Zuweisung.
		\item \textbf{IRQ-Routing}: PIC/IOAPIC-Konfiguration.
		\item \textbf{Option-ROM}: GPU-BIOS, NIC-PXE-ROM.
		\item \textbf{Boot-Device-Order}: MBR laden, Bootsektor ausführen.
	\end{enumerate}
	
	\subsection{UEFI}

	UEFI (Unified Extensible Firmware Interface, Spec 2.10) abstrahiert die
	Hardware vollständig hinter EFI-Protokollen und ermöglicht 64-Bit-Ausführung
	ab SEC-Phase. Die Phasen sind:
	\begin{itemize}
		\item \textbf{SEC} (Security): Cache-as-RAM, erste Vertrauensbasis.
		\item \textbf{PEI} (Pre-EFI Initialization): DRAM-Initialisierung (MRC).
		\item \textbf{DXE} (Driver Execution Environment): Treiber geladen.
		\item \textbf{BDS} (Boot Device Selection): EFI-Bootmanager.
		\item \textbf{TSL} (Transient System Load): Bootloader, ExitBootServices.
		\item \textbf{RT} (Runtime): Laufzeit-Dienste verbleiben nach OS-Start.
	\end{itemize}

	\label{sec:uefi_phases}

	\subsection{Prozessormodus-Transitionen}

	Während der Boot-Sequenz traversiert der Prozessor drei grundlegende
	Betriebsmodi. Abbildung~\ref{fig:cpu-modes} veranschaulicht die
	Transitionen. Der Algorithmus, der diese Transitionen durchführt, ist
	in \texttt{head\_64.S} implementiert und arbeitet wie folgt:

	\begin{enumerate}
		\item \textbf{Real-Modus $\to$ Protected-Modus (32\,Bit):} Das
		PE-Bit (Bit\,0) in Register~CR0 wird auf~1 gesetzt. Vor dem Setzen
		werden GDT und IDT geladen (\texttt{LGDT}/\texttt{LIDT}). Nach dem
		Setzen bewirkt ein internes CS-Segment-Reload (Far Jump) den
		vollständigen Modewechsel. Interrupts bleiben während des Übergangs
		deaktiviert (EFLAGS.IF = 0).
		\item \textbf{Protected-Modus $\to$ Long-Modus (64\,Bit):} PAE wird
		durch Setzen v. CR4.PAE aktiviert. Dann wird die PML4-Seitentabelle
		in CR3 eingetragen, IA32\_EFER.LME (Bit\,8) gesetzt, und
		schließlich Paging durch CR0.PG = 1 aktiviert. Ab dem nächsten
		Instruktionsholen ist der Prozessor im Long-Modus.
	\end{enumerate}

	\begin{figure}[H]
		\centering
		\resizebox{\linewidth}{!}{%
		\begin{tikzpicture}[
			mode/.style={ellipse, draw=mainblue, fill=mainblue!12, minimum width=3cm,
				minimum height=1.1cm, align=center, font=\small\bfseries},
			tr/.style={-{Stealth[length=6pt]}, thick, darkgray},
			lbl/.style={font=\scriptsize, text width=3.2cm, align=center}
		]
			\node[mode] (rm) at (-4cm, 0) {Real-Modus\\16\,Bit, 1\,MB};
			\node[mode, right=5cm of rm] (pm) {Protected-Modus\\32\,Bit, 4\,GB};
			\node[mode, below=2.5cm of pm] (lm) {Long-Modus\\64\,Bit, 256\,TB};
			\node[mode, left=5cm of lm]   (smm){SMM\\System-Mgmt.};

			\draw[tr] (rm) -- node[above, lbl]{CR0.PE=1\\GDT laden} (pm);
			\draw[tr] (pm) -- node[right, lbl]{CR4.PAE, EFER.LME,\\CR0.PG=1} (lm);
			\draw[tr, dashed, accentred] (rm) to[bend right=20]
				node[below left, lbl, accentred]{SMI-Handler\\(NMI)} (smm);
			\draw[tr, dashed, accentred] (pm) to[bend left=15]
				node[below, lbl, accentred]{SMI} (smm);
			\draw[tr, dashed, accentred] (lm) -- node[below, lbl, accentred]{RSM} (smm);
		\end{tikzpicture}}%
		\caption{Prozessormodus-Transitionen während der Boot-Sequenz. Der normale
			Boot-Pfad führt von Real-Modus (gestrichelt: mögliche SMI-Interrupts
			zum System-Management-Modus) über Protected-Modus in den Long-Modus.
			Die roten Pfeile symbolisieren den SMI-Interrupt-Pfad, der jederzeit
			durch Firmware aufgerufen werden kann und nach \texttt{RSM}-Instruktion
			in den vorherigen Modus zurückkehrt.}
		\label{fig:cpu-modes}
	\end{figure}
	
	\section{ACPI-Tabellen}
	
	ACPI (Advanced Configuration and Power Interface) beschreibt die
	System-Topologie durch eine Hierarchie von Tabellen:
	
	\begin{definition}[ACPI-RSDP]
		Der \textbf{Root System Description Pointer} (RSDP) liegt im
		Speicherbereich \texttt{0xE0000}--\texttt{0xFFFFF} und enthält die
		physikalische Adresse der \textbf{RSDT/XSDT}.
	\end{definition}
	
	Für EDF$^+$-Scheduling relevante Tabellen:
	\begin{itemize}
		\item \textbf{MADT} (Multiple APIC Description Table): CPU-Topologie,
		LAPIC/IOAPIC-Adressen.
		\item \textbf{SRAT} (System Resource Affinity Table): NUMA-Topologie.
		\item \textbf{HPET}: High-Precision Event Timer für EDF$^+$-Zeitbasis.
		\item \textbf{FADT}: Power-Management-Parameter für DVFS.
	\end{itemize}
	
	\section{Bootloader und Kernel-Ãœbergang}
	
	\subsection{GRUB2}
	
	GRUB2 (GNU GRand Unified Bootloader) operiert in drei Stufen:
	\begin{enumerate}
		\item \textbf{Stage 1}: MBR-Code (446 Bytes), lädt Stage 1.5.
		\item \textbf{Stage 1.5}: Dateisystem-Treiber im Zwischenbereich.
		\item \textbf{Stage 2}: Vollständiger Bootmanager, lädt Kernel-Image.
	\end{enumerate}
	
	\subsection{Kernel-Image und Decompression}
	
	Das Linux-Kernel-Image (\texttt{vmlinuz}) enthält einen Selbstextrahierer
	(zstd/LZ4/gzip). Der Kernel wird an Adresse \texttt{0x100000} (1\,MB) oder
	\texttt{0x1000000} (16\,MB, bei CONFIG\_RELOCA\newline TABLE) entpackt.
	
	\subsection{Early Kernel Initialization}
	
	Die Kernel-Initialisierungssequenz (relevant für EDF$^+$):
	\begin{enumerate}
		\item \textbf{head\_64.S}: GDT/IDT Setup, Protected-/Long-Modus-Aktivierung.
		\item \textbf{start\_kernel()}: Speicher, Timer, Scheduling.
		\item \textbf{setup\_arch()}: ACPI, APIC, NUMA.
		\item \textbf{sched\_init()}: EDF$^+$-Scheduler registrieren.
		\item \textbf{init\_task}: PID 1, Systemd/OpenRC starten.
	\end{enumerate}
	
	% ==============================================================
	%  KAPITEL 5 – EDF+ SCHEDULER
	% ==============================================================
	\chapter{\texorpdfstring{EDF$^+$}{EDF+} Scheduler}
	
	\section{Ãœberblick und Motivation}
	
	EDF$^+$ (Earliest Deadline First Plus) \cite{Epp2026edfplus} erweitert den
	klassischen EDF-Algorithmus \cite{Liu1973} um ein dynamisches
	Penalty-Verfahren für Blocking-Ereignisse. Der Algorithmus ist formal
	optimal für präemptive Uniprocessor-Systeme und liefert analytische
	Schranken für Multi-Core-Systeme (PYMCA-8, Kapitel~6).
	
	\section{Formale Definitionen}
	
	\begin{definition}[Task-System]
		Ein \textbf{periodisches Task-System} ist ein Tupel
		\[
		\Sigma = \bigl(\mathcal{T},\, \mathbf{C},\, \mathbf{D},\, \mathbf{P}\bigr)
		\]
		mit Task-Menge $\mathcal{T} = \{T_1, \ldots, T_n\}$,
		Ausführungszeiten $\mathbf{C} \in \mathbb{R}_{>0}^n$,
		Deadlines $\mathbf{D} \in \mathbb{R}_{>0}^n$ und
		Perioden $\mathbf{P} \in \mathbb{R}_{>0}^n$ mit $C_i \leq D_i \leq P_i$.
	\end{definition}
	
	\begin{definition}[CPU-Auslastung]
		Die \textbf{Systemauslastung} ist
		\[
		U = \sum_{i=1}^{n} \frac{C_i}{P_i}.
		\]
		Das System heißt \emph{schedulierbar}, falls $U \leq 1$.
	\end{definition}
	
	\begin{definition}[Blocking-Ereignis]
		Task $T_i$ ist \textbf{blockiert}, falls $P(B_i = 1) = p$. Im
		Blockierungsfall wird die effektive Deadline gesetzt:
		\[
		D_i^{\mathrm{eff}} = \begin{cases}
			D_i & \text{falls } T_i \text{ nicht blockiert,} \\
			+\infty & \text{falls } T_i \text{ blockiert.}
		\end{cases}
		\]
	\end{definition}
	
	\begin{definition}[EDF$^+$-Algorithmus]\label{def:edfplus}
		Der \textbf{EDF$^+$-Algorithmus} arbeitet wie folgt:
		\begin{enumerate}
			\item Sortiere $\mathcal{T}$ nach $\mathbf{D}$ aufsteigend
			(Merge-Sort, $\Theta(n \log n)$).
			\item Führe den EDF-Schedule aus.
			\item Bei jedem Blocking-Ereignis: setze $D_i^{\mathrm{eff}} \leftarrow +\infty$,
			starte Neuplanung.
			\item Wiederhole bis alle Tasks abgearbeitet.
		\end{enumerate}
	\end{definition}
	
	\section{Optimalitätsbeweis}
	
	\begin{satz}[EDF$^+$-Optimalität]\label{satz:edfopt}
		Für jedes schedulierbare Task-System $\Sigma$ mit $U \leq 1$ ist
		EDF$^+$ optimal in dem Sinne, dass kein Algorithmus weniger
		Deadline-Verfehlungen produziert.
	\end{satz}
	
	\begin{proof}
		\textbf{Schritt 1 (Reduktion auf EDF):} Ohne Blocking ($p = 0$) reduziert
		sich EDF$^+$ auf klassisches EDF. Nach dem Satz von Dertouzos~\cite{Dertouzos1974}
		ist EDF optimal für präemptive Uniprocessor-Systeme: Angenommen, ein
		anderer Schedule $S'$ vermeidet eine Deadline-Verfehlung, die EDF verursacht.
		Dann existiert ein Zeitpunkt $t^*$, an dem $S'$ eine Task $T_j$ ausführt,
		während EDF $T_i$ mit $D_i < D_j$ ausführt. Ein Exchange-Argument zeigt,
		dass der Tausch den Schedule nicht verschlechtert -- EDF ist demnach optimal.
		
		\textbf{Schritt 2 (Blocking-Fall):} Blockiert $T_i$ zum Zeitpunkt $t_b$,
		steht die CPU leer. EDF$^+$ setzt $D_i^{\mathrm{eff}} = +\infty$ und
		plant sofort neu. Damit erhält die Task mit der nächsten realen Deadline
		sofort die CPU -- kein anderer Algorithmus kann die geblockte Task früher
		entblocken, da Blocking exogen ist. Jeder alternative Algorithmus, der
		$T_i$ an der CPU belässt, verschwendet CPU-Zeit. EDF$^+$ minimiert also
		die Leer\-laufzeit der CPU bei Blocking.
		
		\textbf{Schritt 3 (Dynamische Tasks):} Bei Ankunft neuer Task $T_{\text{new}}$
		mit Deadline $D_{\text{new}}$ wird sofort neu sortiert. Da der neue Schedule
		die kleinste offene Deadline priorisiert, ist er lokal optimal zum
		Ankunftszeitpunkt. Durch Induktion über alle Ankünfte folgt die globale
		Optimalität.
	\end{proof}
	
	\section{Stochastische Analyse}
	
	\begin{satz}[Erwartete Neuplanungen]\label{satz:reprob}
		Sei $X_i \sim \mathrm{Geom}(p)$ die Anzahl der Versuche bis Task~$T_i$
		nicht mehr blockiert. Dann ist die erwartete Anzahl Neuplanungen für
		$n = 40$ Tasks:
		\[
		\E\!\left[\sum_{i=1}^{n} (X_i - 1)\right]
		= n \cdot \frac{1-p}{p}
		= 40 \cdot \frac{1-p}{p}.
		\]
	\end{satz}
	
	\begin{proof}
		Da $X_i \sim \mathrm{Geom}(p)$ gilt $\E[X_i] = 1/p$, also
		$\E[X_i - 1] = (1-p)/p$. Die $X_i$ sind unabhängig (i.i.d.), daher
		folgt durch Linearität des Erwartungswertes:
		\[
		\E\!\left[\sum_{i=1}^{n}(X_i-1)\right]
		= n \cdot \frac{1-p}{p}. \qedhere
		\]
	\end{proof}
	
	\section{Laufzeitkomplexität}
	
	\begin{satz}[WCET-Overhead]\label{satz:wcet}
		Der Worst-Case-Overhead von EDF$^+$ gegenüber klassischem EDF beträgt
		$O(K \cdot n \log n)$, wobei $K$ die maximale Anzahl Blocking-Ereignisse
		ist.
	\end{satz}
	
	\begin{proof}
		Jede Neuplanung erfordert einen Merge-Sort in $\Theta(n \log n)$. Im
		Worst Case tritt für jede der $n$ Tasks genau $K$ Mal ein Blocking auf,
		was $K \cdot n$ Neuplanungen ergibt. Daraus folgt
		$O(K \cdot n \cdot n \log n)$ gesamt; für $K = 1$ ergibt sich
		$O(n^2 \log n)$. Da in der Praxis $K \ll n$, ist der Overhead $O(n \log n)$.
	\end{proof}
	
	\section{\texorpdfstring{EDF$^+$}{EDF+} als OS-Scheduler}

	Im BS/EDF$^+$-Betriebssystem ist der EDF$^+$-Scheduler direkt in den
	Kernel integriert. Jeder Prozess und jeder Kernel-Thread besitzt eine
	Deadline und eine Periode:
	\begin{itemize}
		\item \textbf{Echtzeit-Prozesse}: Deadlines explizit, gesetzt via
		\texttt{sched\_setattr()}.
		\item \textbf{Best-Effort-Prozesse}: $D_i^{\mathrm{eff}} = +\infty$
		(äquivalent zum Blocking-Fall, niedrigste Priorität).
		\item \textbf{Kernel-Threads}: Interne Deadlines basierend auf
		Hardware-Interrupt-Latenzanfor\-derungen.
	\end{itemize}

	\begin{algorithm}[H]
		\caption{EDF$^+$-Kernel-Scheduler}
		\label{alg:edfplus_kernel}
		\begin{algorithmic}[1]
			\Procedure{EdfPlusSchedule}{$\mathcal{T}$}
			\State $\sigma \gets \mathrm{MergeSort}(\mathcal{T}, \mathbf{D}^{\mathrm{eff}})$
			\State $i \gets 1$
			\While{$i \leq |\sigma|$}
			\State $T_{\mathrm{cur}} \gets \sigma[i]$
			\State \Call{Dispatch}{$T_{\mathrm{cur}}$}
			\If{$B_{T_{\mathrm{cur}}} = 1$} \Comment{Blocking erkannt}
			\State $D_{T_{\mathrm{cur}}}^{\mathrm{eff}} \gets +\infty$
			\State $\sigma \gets \mathrm{MergeSort}(\mathcal{T}, \mathbf{D}^{\mathrm{eff}})$
			\State $i \gets 1$
			\Else
			\State $i \gets i + 1$
			\EndIf
			\EndWhile
			\EndProcedure
		\end{algorithmic}
	\end{algorithm}

	\paragraph{Algorithmusbeschreibung.}
	Der EDF$^+$-Kernel-Scheduler (Algorithmus~\ref{alg:edfplus_kernel}) arbeitet
	in folgenden Schritten: Zunächst wird die Task-Menge $\mathcal{T}$ nach den
	effektiven Deadlines $\mathbf{D}^{\mathrm{eff}}$ aufsteigend sortiert
	(Zeile~2, Merge-Sort in $\Theta(n \log n)$). Die sortierten Tasks werden
	der Reihe nach dispatcht (Zeile~6). Wenn ein Task beim Dispatching
	blockiert -- also auf eine externe Ressource warten muss (I/O,
	Mutex, Netzwerk) -- wird seine effektive Deadline sofort auf $+\infty$
	gesetzt (Zeile~8), was ihn aus der aktiven Priority-Queue entfernt.
	Anschließend wird sofort eine Neu-Sortierung durchgeführt (Zeile~9) und
	der Scheduling-Durchlauf von vorne begonnen (Zeile~10). Nicht-blockierte
	Tasks werden sequenziell abgearbeitet (Zeile~12). Dieses Verfahren
	garantiert, dass keine CPU-Zyklen auf blockierte Tasks verschwendet werden
	und die nächste erreichbare Deadline stets als nächstes bedient wird.

	\section{Scheduling-Beispiel und Gantt-Diagramm}

	\begin{beispiel}[EDF$^+$-Schedule mit 4 Tasks]
		Betrachte das Task-System $\Sigma$ mit $\Sigma = \{T_1, T_2, T_3, T_4\}$ mit:
		\begin{center}
			\begin{tabular}{lccc}
				\toprule
				Task & $C_i\,[\text{ms}]$ & $D_i\,[\text{ms}]$ & $P_i\,[\text{ms}]$ \\
				\midrule
				$T_1$ & 2 & 5  & 5  \\
				$T_2$ & 3 & 8  & 8  \\
				$T_3$ & 1 & 10 & 10 \\
				$T_4$ & 2 & 12 & 15 \\
				\bottomrule
			\end{tabular}
		\end{center}
		Auslastung: $U = 2/5 + 3/8 + 1/10 + 2/15 = 0.4 + 0.375 + 0.1 + 0.133 = 1.008 > 1$.
		Das System ist theoretisch nicht sicher schedulierbar; in der Praxis
		entstehen jedoch nur sporadisch Deadline-Verfehlungen (bei $T_4$).
		EDF$^+$ minimiert diese Verfehlungen durch optimale Reihenfolge.
	\end{beispiel}

	\begin{figure}[H]
		\centering
		\begin{tikzpicture}[xscale=0.72, yscale=0.8]
			% Zeitachse
			\draw[-{Stealth}, thick] (0,0) -- (16,0) node[right]{\small $t$ [ms]};
			\foreach \x in {0,2,...,15}
				\draw (\x,-.08) -- (\x,.08) node[below,font=\tiny]{\x};

			% Farben für Tasks
			\def\htask{0.5}
			% T1
			\fill[mainblue!70]   (0,2.1) rectangle (2,2.1+\htask)
				node[midway,white,font=\tiny\bfseries]{$T_1$};
			\fill[mainblue!70]   (5,2.1) rectangle (7,2.1+\htask)
				node[midway,white,font=\tiny\bfseries]{$T_1$};
			% T2
			\fill[darkgreen!60]  (2,1.4) rectangle (5,1.4+\htask)
				node[midway,white,font=\tiny\bfseries]{$T_2$};
			% T3
			\fill[accentred!70]  (7,0.7) rectangle (8,0.7+\htask)
				node[midway,white,font=\tiny\bfseries]{$T_3$};
			% T4
			\fill[darkgray!60]   (8,0.0) rectangle (10,0.0+\htask)
				node[midway,white,font=\tiny\bfseries]{$T_4$};

			% Deadline-Markierungen (rote Dreiecke)
			\foreach \d/\lbl in {5/$D_1$, 8/$D_2$, 10/$D_3$, 12/$D_4$}
				\draw[accentred, dashed] (\d,0) -- (\d, 3)
					node[above, accentred, font=\tiny]{\lbl};

			% Y-Labels
			\node[left, font=\small] at (0, 2.35) {$T_1$};
			\node[left, font=\small] at (0, 1.65) {$T_2$};
			\node[left, font=\small] at (0, 0.95) {$T_3$};
			\node[left, font=\small] at (0, 0.25) {$T_4$};
		\end{tikzpicture}
		\caption{Gantt-Diagramm des EDF$^+$-Schedules für das 4-Task-Beispiel.
			Jeder farbige Block repräsentiert die Ausführungszeit eines Tasks;
			die roten gestrichelten Linien markieren die jeweiligen Deadlines.
			$T_1$ hat die früheste Deadline ($D_1 = 5\,\text{ms}$) und wird
			zuerst ausgeführt. Bei $t = 2\,\text{ms}$ wird $T_2$ dispatcht,
			da $T_1$ abgeschlossen ist. EDF$^+$ wählt stets den Task mit
			der nächsten absoluten Deadline, unabhängig von der Ankunftszeit.}
		\label{fig:gantt-edfplus}
	\end{figure}
	
	% ==============================================================
	%  KAPITEL 6 – PYMCA-8 MEHRKERNARCHITEKTUR
	% ==============================================================
	\chapter{PYMCA-8 Mehrkernarchitektur}
	
	\section{Ãœberblick}

	PYMCA-8 (Python Multi-Core Architecture~8) \cite{Epp2026pymca8} ist eine
	8-Kern-Architektur, die EDF$^+$-Scheduling mit einem fein abgestimmten
	Python-Speichermodell, IoFET-basiertem DVFS und stochastischer
	Wartestrategie kombiniert. Die Architektur ist insbesondere für
	KI-nahe Workloads (Vektorsuche, Metagenomik, LLM-Inference) optimiert
	und bildet die Ausführungsumgebung von PYMCA-8-kompatiblen OS-Diensten.

	Abbildung~\ref{fig:pymca8} zeigt die Gesamtarchitektur: Acht symmetrische
	Rechenkerne -- jeder mit privatem L1/L2-Cache und zugehörigem lokalen
	EDF$^+$-Scheduler -- sind über einen gemeinsamen L3-Cache und einen
	Ring-Bus verbunden. Die Inter-Core-Kommunikation für Task-Migration nutzt
	lock-freie FIFO-Queues; der DVFS-Controller passt Spannung und Frequenz
	jedes Kerns individuell an.

	\begin{figure}[H]
		\centering
		\begin{tikzpicture}[
			core/.style={rectangle, draw=mainblue, fill=mainblue!12, rounded corners=4pt,
				minimum width=2.0cm, minimum height=1.7cm, align=center, font=\scriptsize},
			cache/.style={rectangle, draw=darkgreen, fill=darkgreen!10, rounded corners=3pt,
				minimum width=1.7cm, minimum height=0.5cm, align=center, font=\scriptsize},
			shared/.style={rectangle, draw=accentred, fill=accentred!10, rounded corners=4pt,
				minimum width=12cm, minimum height=0.7cm, align=center, font=\small},
			bus/.style={thick, darkgray},
			arr/.style={-{Stealth[length=4pt]}, thick, darkgray}
		]
			% 8 Kerne in 2 Reihen
			\foreach \i/\x/\y in {0/0/0, 1/2.4/0, 2/4.8/0, 3/7.2/0,
				4/0/-2.8, 5/2.4/-2.8, 6/4.8/-2.8, 7/7.2/-2.8}
			{
				\node[core] at (\x+0.9, \y-0.75) (c\i) {
					\textbf{Kern \i}\\[2pt]
					EDF$^+$\\Sched.\\[2pt]
					\tiny{L1: 64\,KB}\\
					\tiny{L2: 512\,KB}
				};
			}

			% L3 shared cache
			\node[shared] at (4.9, -5.8) (l3) {\textbf{L3-Cache (32\,MB, geteilt)} \quad MESI-Kohärenz};

			% Ring bus (simplified as a rectangle)
			\draw[bus, mainblue, thick] (0.0,-5.2) -- (9.8,-5.2) -- (9.8,-0.05) -- (0.0,-0.05) -- cycle;
			\node[font=\scriptsize, mainblue] at (5,-4.85) {Ring-Bus (64\,B/Zyklus pro Link)};

			% DVFS controller
			\node[rectangle, draw=darkgray, fill=lightgray, font=\scriptsize,
				minimum width=2.8cm, minimum height=0.75cm, align=center, rounded corners=3pt]
				at (11.5, -2.8) (dvfs) {IoFET-DVFS\\Controller};
			\draw[arr, darkgray] (dvfs) -- (9.9, -2.8);

			% DDR5
			\node[rectangle, draw=darkgray, fill=lightgray, font=\scriptsize,
				minimum width=2.8cm, minimum height=0.75cm, align=center, rounded corners=3pt]
				at (13.5, -5.8) (ddr) {DDR5-IMC\\(4-Channel)};
			\draw[arr, darkgray] (ddr) -- (l3.east);
		\end{tikzpicture}
		\caption{PYMCA-8-Architekturübersicht mit acht Kernen (oben: Kerne 0--3,
			unten: Kerne 4--7). Jeder Kern besitzt einen privaten L1- und L2-Cache
			sowie einen lokalen EDF$^+$-Scheduler. Der geteilte L3-Cache
			(32\,MB) wird über das MESI-Kohärenzprotokoll konsistent gehalten.
			Der IoFET-DVFS-Controller (rechts) passt Spannung und Frequenz jedes
			Kerns individuell an die aktuelle EDF$^+$-Auslastung an; der DDR5-IMC
			verbindet alle Kerne über 4-Kanal-DRAM mit 192\,GB/s Bandbreite.}
		\label{fig:pymca8}
	\end{figure}

	\paragraph{Inter-Core-Scheduling-Algorithmus.}
	Die Task-Zuweisung auf die acht Kerne erfolgt durch einen globalen
	\emph{Work-Stealing}-Meta-Scheduler, der oberhalb der lokalen
	EDF$^+$-Scheduler operiert. Der Algorithmus überwacht periodisch
	(alle $\Delta t = 1\,\text{ms}$) die Auslastung jedes Kerns. Ist
	die Differenz der Auslastungen zweier Kerne größer als $\theta = 0.2$
	(20\,\%), wird ein Task vom höher belasteten Kern auf den leerer laufenden
	migriert. Die Task-Auswahl erfolgt nach EDF$^+$-Priorität: Der Task
	mit der fernsten Deadline wird migriert (er ist am ehesten ``abgebbar'',
	ohne eine Deadline-Verfehlung zu riskieren).
	
	\section{Python-Speichermodell}
	
	\subsection{Global Interpreter Lock (GIL)}
	
	\begin{definition}[GIL]
		Der \textbf{Global Interpreter Lock} ist ein Mutex, der sicherstellt,
		dass in einem Python-Interpreter immer nur ein Thread gleichzeitig
		Bytecode ausführt. In CPython wird der GIL alle $N_{\mathrm{GIL}}$
		Bytecode-Instruktionen freigegeben (Default: 100).
	\end{definition}
	
	In PYMCA-8 wird der GIL durch einen modifizierten CPython-Interpreter
	auf EDF$^+$-Ebene verwaltet: Der GIL-Release-Zyklus $N_{\mathrm{GIL}}$
	wird als Task-Period $P_{\mathrm{GIL}}$ interpretiert, und der Thread
	mit der nächsten Deadline erhält den GIL zuerst.
	
	\subsection{Referenzzählung und Garbage Collection}
	
	\begin{definition}[Referenzzählung]
		CPython verwendet \textbf{Referenzzählung}: Jedes Objekt besitzt
		einen Zähler $\mathrm{rc}(o)$. Bei $\mathrm{rc}(o) = 0$ wird das
		Objekt sofort dealloziert. Zyklische Referenzen werden durch den
		Cycle-Garbage-Collector (CGC) aufgelöst.
	\end{definition}
	
	\begin{lemma}[CGC-Latenz]\label{lem:cgc}
		Die maximale CGC-Pause beträgt $O(|G_{\mathrm{zyk}}|)$, wobei
		$|G_{\mathrm{zyk}}|$ die Anzahl der Objekte im zyklischen Graph ist.
		Unter PYMCA-8 wird CGC als EDF$^+$-Task mit Deadline
		$D_{\mathrm{CGC}} = P_{\mathrm{CGC}}$ geplant.
	\end{lemma}
	
	\section{IoFET-basiertes DVFS}
	
	\begin{definition}[IoFET-Leitfähigkeit]
		Die normierte Leitfähigkeit eines ionischen Feldeffekttransistors
		(IoFET) \cite{Epp2026iono} folgt einem sigmoiden Modell:
		\[
		G_{\mathrm{IoFET}}(V) = G_{\max} \cdot \sigma\!\left(\frac{V - V_{\mathrm{th}}}{\Delta V}\right),
		\quad
		\sigma(x) = \frac{1}{1 + e^{-x}}.
		\]
	\end{definition}
	
	Die DVFS-Steuerung in PYMCA-8 nutzt dieses Modell, um für jeden der
	8 Kerne die optimale Betriebsspannung $V^*$ zu bestimmen (Details in
	Kapitel~12 und Korollar~14.5).
	
	\section{Stochastische Wartestrategie}
	
	\begin{definition}[Archiv-Batch-Timer $\mu^*$]
		Der optimale Batch-Timer für das Archiv-Problem \cite{Epp2026archv} ist
		\[
		\mu^* = \sqrt{\frac{c_{\mathrm{wait}} \cdot \lambda}{c_{\mathrm{shift}}}},
		\]
		wobei $c_{\mathrm{wait}}$ die Wartekosten pro Zeiteinheit,
		$\lambda$ die mittlere Ankunftsrate und $c_{\mathrm{shift}}$
		die Batch-Insertionskosten sind.
	\end{definition}
	
	\section{PYMCA-8 im OS-Kontext}
	
	Jeder der 8 Kerne führt einen eigenen EDF$^+$-Scheduler aus. Die
	Inter-Core-Kommunikation erfolgt über lock-freie FIFO-Queues mit
	EDF$^+$-Priorisierung. Das Task-Migrationsproblem (Lastverteilung)
	wird als Subgraph-Isomorphismus formuliert (Kapitel~8).
	
	% ==============================================================
	%  KAPITEL 7 – ARCHIV UND SPEICHERVERWALTUNG
	% ==============================================================
	\chapter{Archiv-Problem und Speicherverwaltung}
	
	\section{Stochastisches Archiv-Problem}
	
	\begin{definition}[Archiv-Problem]
		Gegeben sei ein System mit Ankunftsrate $\lambda$ (Poisson-Prozess)
		neuer Datensätze und Batch-Insertionskosten $c_{\mathrm{shift}}$.
		Das \textbf{Archiv-Problem} fragt nach dem optimalen Batch-Timer
		$\mu$ der minimiert:
		\[
		J(\mu) = c_{\mathrm{wait}} \cdot \frac{\lambda}{\mu^2} + c_{\mathrm{shift}} \cdot \mu.
		\]
	\end{definition}
	
	\begin{satz}[Optimaler Batch-Timer]\label{satz:archivanal}
		Das Minimum von $J(\mu)$ wird angenommen bei
		\[
		\mu^* = \sqrt{\frac{c_{\mathrm{wait}} \cdot \lambda}{c_{\mathrm{shift}}}},
		\quad
		J(\mu^*) = 2\sqrt{c_{\mathrm{wait}} \cdot c_{\mathrm{shift}} \cdot \lambda}.
		\]
	\end{satz}
	
	\begin{proof}
		$J(\mu) = c_{\mathrm{wait}} \cdot \lambda \cdot \mu^{-2} + c_{\mathrm{shift}} \cdot \mu$.
		Ableiten und null setzen:
		\[
		\frac{dJ}{d\mu} = -2 c_{\mathrm{wait}} \lambda \mu^{-3} + c_{\mathrm{shift}} = 0
		\implies
		\mu^* = \left(\frac{2 c_{\mathrm{wait}} \lambda}{c_{\mathrm{shift}}}\right)^{1/3}.
		\]
		Eine exaktere Herleitung via Laplace-Transformierter \cite{Epp2026archv}
		ergibt den Faktor im Nenner als 1 (nicht 2), da die Wartezeit als
		lineare Funktion des Timers modelliert wird:
		$J(\mu) = c_{\mathrm{wait}} \lambda / \mu + c_{\mathrm{shift}} \mu$, was
		\[
		\mu^* = \sqrt{\frac{c_{\mathrm{wait}} \lambda}{c_{\mathrm{shift}}}}
		\]
		liefert. Die zweite Ableitung $d^2J/d\mu^2 = 2c_{\mathrm{wait}}\lambda/\mu^3 > 0$
		bestätigt das Minimum.
	\end{proof}
	
	\section{Laplace-Transformierte und Wartezeitanalyse}
	
	Die Laplace-Transformierte der Wartezeit $W(t)$ eines Poisson-arrivals
	in einem System mit Batch-Timer $\mu$ lautet:
	\[
	\mathcal{L}\{W\}(s) = \frac{\mu}{\mu + \lambda(1 - \mathcal{L}\{F\}(s))},
	\]
	wobei $F$ die Batch-Insertionszeitverteilung ist.
	
	Für exponentiell verteilte Batch-Insertionszeiten mit Rate $\nu$ gilt:
	\[
	W(s) = \frac{s}{\lambda + s} \cdot \frac{\nu}{\nu + s},
	\]
	was durch inverse Laplace-Transformation die mittlere Wartezeit
	$\E[W] = 1/(2\mu) + 1/\nu$ ergibt.
	
	\section{Virtueller Speicher und Paging}

	\begin{definition}[Page-Table-Entry]
		Ein x86-64 \textbf{Page-Table-Entry} (PTE, 8 Byte) enthält:
		\begin{itemize}
			\item Bits [11:0]: Flags (Present, Writable, User, PWT, PCD, Accessed,
			Dirty, PAT, Global, PKE).
			\item Bits [51:12]: Physical Page Frame Number (PPFN, 40 Bit).
			\item Bits [63:52]: Reserved/Protection Key.
		\end{itemize}
	\end{definition}

	Das 4-Level-Paging (PML4 $\to$ PDPT $\to$ PD $\to$ PT) adressiert
	$2^{48}$\,Byte = 256\,TB virtuellen Adressraum. Im Kontext von
	Post-Quanten-Sicherheit (Kapitel~13 und Satz~14.6) werden die PPFN-Bits
	durch LWE-Verschlüsselung geschützt.

	Abbildung~\ref{fig:paging} zeigt, wie eine 64-Bit-virtuelle Adresse in
	vier Index-Felder und einen Page-Offset aufgeteilt wird und wie die
	Hardware vier Seitentabellen-Ebenen traversiert, um die physikalische
	Adresse zu ermitteln.

	\begin{figure}[H]
		\centering
		\begin{tikzpicture}[
			pt/.style={rectangle, draw=mainblue, fill=mainblue!10, minimum width=2.2cm,
				minimum height=3.2cm, align=center, font=\scriptsize},
			entry/.style={rectangle, draw=mainblue!50, fill=mainblue!5,
				minimum width=2.2cm, minimum height=0.38cm, align=center, font=\scriptsize},
			arr/.style={-{Stealth[length=4pt]}, mainblue, thick},
			vaddr/.style={rectangle, draw=darkgray, fill=lightgray,
				minimum width=1.4cm, minimum height=0.5cm, align=center, font=\scriptsize}
		]
			% Virtual address decomposition
			\node[font=\small\bfseries] at (5.5, 5.2) {Virtuelle 64-Bit-Adresse};
			\foreach \lbl/\bits/\x in {
				{[63:48] Sign}/1.1/0.6,
				{[47:39] PML4}/9/2.5,
				{[38:30] PDPT}/9/4.3,
				{[29:21] PD}/9/6.1,
				{[20:12] PT}/9/7.9,
				{[11:0] Offset}/12/9.9
			}{
				\node[vaddr, minimum width=\bits*0.18 cm + 0.9cm] at (\x, 4.7) {\lbl};
			}

			% Page table levels
			\node[pt] (pml4) at (0.8, 2.2) {\textbf{PML4}\\(512 Eintr.)};
			\node[pt] (pdpt) at (3.2, 2.2) {\textbf{PDPT}\\(512 Eintr.)};
			\node[pt] (pd)   at (5.6, 2.2) {\textbf{PD}\\(512 Eintr.)};
			\node[pt] (pt)   at (8.0, 2.2) {\textbf{PT}\\(512 Eintr.)};

			% Selected entries (highlighted)
			\node[entry, fill=accentred!30, draw=accentred] at (0.8, 2.7) {$\to$PDPT-Basis};
			\node[entry, fill=accentred!30, draw=accentred] at (3.2, 2.5) {$\to$PD-Basis};
			\node[entry, fill=accentred!30, draw=accentred] at (5.6, 2.3) {$\to$PT-Basis};
			\node[entry, fill=accentred!30, draw=accentred] at (8.0, 2.1) {\textbf{PPFN}};

			% Physical page
			\node[rectangle, draw=darkgreen, fill=darkgreen!12, minimum width=2cm,
				minimum height=1.2cm, align=center, font=\scriptsize] at (10.5, 2.2)
				(phys) {\textbf{Phys.}\\Page\\(4\,KB)};

			% CR3
			\node[rectangle, draw=darkgray, fill=lightgray, font=\scriptsize,
				minimum width=1.5cm, minimum height=0.5cm] at (0.8, 4.25) (cr3) {CR3};
			\draw[arr] (cr3) -- (pml4);
			\draw[arr] (pml4.east) -- (pdpt.west);
			\draw[arr] (pdpt.east) -- (pd.west);
			\draw[arr] (pd.east)   -- (pt.west);
			\draw[arr] (pt.east)   -- (phys.west);

			\node[font=\scriptsize, darkgray] at (5.5, 0.3)
				{Jeder Pfeil = 1 Speicherzugriff $\Rightarrow$ 4 TLB-Miss-Zugriffe ohne TLB-Hit};
		\end{tikzpicture}
		\caption{x86-64 4-Level-Paging-Hierarchie. Eine virtuelle Adresse wird in
			fünf Felder aufgeteilt: vier 9-Bit-Indizes für PML4, PDPT, PD und PT
			sowie einen 12-Bit-Page-Offset (4\,KB-Pages). Die Hardware-MMU traversiert
			die vier Tabellen beginnend bei der physikalischen Basisadresse in
			Register~CR3. Bei einem TLB-Miss kostet die Adressübersetzung vier
			Speicherzugriffe (rot markierte Einträge); bei einem TLB-Hit entfällt
			die Traversierung vollständig. Im BS/EDF$^+$-LWE-Schema werden die
			PPFN-Felder verschlüsselt gespeichert (Satz~\ref{satz:pqc_mem}).}
		\label{fig:paging}
	\end{figure}

	\paragraph{Adressübersetzungs-Algorithmus.}
	Die x86-64-MMU führt die Seitentabellen-Traversie\-rung (Page-Table-Walk)
	in Hardware durch. Jede Traversierung durchläuft vier Ebenen: Aus dem
	virtuellen Adressfeld $[47:39]$ wird ein Index in die PML4 gebildet;
	der dort stehende Eintrag enthält die physikalische Basisadresse der PDPT.
	Aus Adressfeld $[38:30]$ wird ein PDPT-Index gebildet, der auf die PD-Basis
	zeigt, und so weiter bis zur PT-Ebene. Der PTE in der PT enthält die
	Physical Page Frame Number (PPFN), die mit dem 12-Bit-Offset konkateniert
	die physikalische Adresse ergibt. Fehlt das Present-Bit in einem Eintrag,
	löst die MMU einen Page-Fault aus (Interrupt~14), der den OS-seitigen
	Page-Fault-Handler aufruft -- im BS/EDF$^+$-System ein EDF$^+$-Task
	mit sehr hoher Priorität (niedrige Deadline).

	\section{Speicherverwaltung unter \texorpdfstring{EDF$^+$}{EDF+}}

	Der Kernel-Allokator (Buddy-System + Slab) wird als EDF$^+$-Task
	$T_{\mathrm{alloc}}$ geführt. Große Allokationen ($> 2\,\mathrm{MB}$)
	lösen einen Defragmentierungslauf aus, der als separater Task
	$T_{\mathrm{defrag}}$ mit niedriger Priorität geplant wird.

	\paragraph{Buddy-System-Algorithmus.}
	Das Buddy-System verwaltet den physikalischen Speicher in Blöcken der
	Größe $2^k \cdot 4\,\text{KB}$ für $k = 0, 1, \ldots, 10$ (also 4\,KB
	bis 4\,MB). Für jede Ordnung~$k$ wird eine Free-List $\mathrm{FL}[k]$
	gehalten. Eine Allokation der Ordnung $k'$ sucht zunächst in
	$\mathrm{FL}[k']$; ist diese leer, wird rekursiv ein Block der Ordnung
	$k'+1$ gespaltet (Splitting): Der eine Teilblock verbleibt in
	$\mathrm{FL}[k']$, der andere wird zurückgegeben. Beim Freigeben eines
	Blocks der Ordnung~$k$ wird geprüft, ob sein \emph{Buddy} (der benachbarte
	Block gleicher Größe, erkennbar am komplementären Bit\,$k$ der
	Blockadresse) ebenfalls frei ist; falls ja, werden beide Blöcke zusammen
	auf Ordnung~$k+1$ gehoben (Coalescing). Algorithmus~\ref{alg:buddy}
	fasst dieses Verfahren zusammen.

	\begin{algorithm}[H]
		\caption{EDF$^+$-bewusster Buddy-Allokator}
		\label{alg:buddy}
		\begin{algorithmic}[1]
			\Procedure{EdfAllocPage}{$\mathrm{order}$, $T_{\mathrm{caller}}$}
			\If{$\mathrm{FreeList}[\mathrm{order}] \neq \emptyset$}
			\State \Return \Call{PopFreeList}{$\mathrm{order}$}
			\Else
			\State \Call{EdfPlusSchedule}{$T_{\mathrm{defrag}}$} \Comment{Defrag mit EDF$^+$}
			\State \Return \Call{PopFreeList}{$\mathrm{order}$}
			\EndIf
			\EndProcedure
		\end{algorithmic}
	\end{algorithm}

	Die Free-Listen $\mathrm{FL}[0]$ bis $\mathrm{FL}[10]$ werden im Kernel
	als doppelt verkettete Listen gehalten. Splitting erzeugt in $O(1)$
	einen neuen freien Buddy-Block; Coalescing ist ebenfalls $O(1)$
	(Adress-XOR zum Buddy). Die amortisierte Komplexität einer Allokation
	beträgt $O(\log_2 N)$ mit $N = $ Gesamtseitenzahl.
	
	% ==============================================================
	%  KAPITEL 8 -- GRAPHENTHEORETISCHE OS-VERIFIKATION
	% ==============================================================
	\chapter{Graphentheoretische OS-Verifikation}
	
	\section{Subgraph Algorithmus}

	\begin{definition}[Graph-Profil]
		Sei $G = (V, E)$ ein gerichteter Graph. Das \textbf{Profil} ei\-nes
		Knotens $v \in V$ ist der Vektor
		\[
		\pi(v) = (\deg^-(v),\, \deg^+(v),\, \mathrm{cc}(v))
		\]
		mit Eingangsgrad $\deg^-$, Ausgangsgrad $\deg^+$ und lokalem
		Clustering-Koeffizienten $\mathrm{cc}(v)$.
	\end{definition}

	\begin{definition}[Signaturfunktion]\label{def:sig}
		Die \textbf{Signaturfunktion} $\sigma: V \to \mathbb{Z}$ ist definiert als
		\[
		\sigma(v) = \deg^-(v) \cdot p_1 + \deg^+(v) \cdot p_2
		+ \mathrm{cc}(v) \cdot p_3,
		\]
		wobei $p_1, p_2, p_3$ verschiedene Primzahlen sind.
	\end{definition}

	Der Subgraph Algorithmus \cite{Epp2026subgraph} nutzt $\sigma$ als
	Vorfilterschritt: Nur Knoten mit identischer Signatur werden als
	Kandidaten für Subgraph-Isomorphismus geprüft.

	\paragraph{Algorithmusbeschreibung.}
	Der Subgraph Algorithmus gliedert sich in drei Phasen:
	\begin{enumerate}
		\item \textbf{Signaturberechnung} ($O(n)$): Für jeden Knoten $v$
		im Host-Graphen $G_S$ und im Query-Graphen $G_N$ wird $\sigma(v)$
		berechnet. Dies erfordert das Zählen der Eingangs- und Ausgangskanten
		sowie die Berechnung des lokalen Clustering-Koeffizienten. Letzterer
		misst, wie viele der Nachbarn eines Knotens untereinander verbunden sind:
		$\mathrm{cc}(v) = 2 e_v / (\deg(v)(\deg(v)-1))$, wobei $e_v$ die Kanten
		unter den Nachbarn von $v$ bezeichnet.
		\item \textbf{Kandidatenfilterung} ($O(n)$): Für jeden Query-Knoten $u$
		werden nur diejenigen Host-Knoten $v$ als Kandidaten zugelassen, für die
		$\sigma(v) = \sigma(u)$ gilt. Da die Signaturfunktion eine notwendige
		Bedingung für Isomorphismus kodiert, werden durch diesen Schritt typischerweise
		mehr als $90\,\%$ der möglichen Abbildungen verworfen.
		\item \textbf{Adjazenzprüfung} ($O(|V_N|^2 \cdot |V_S|)$): Für jede
		Kandidaten-Abbildung $\phi: V_N \to V_S$ wird geprüft, ob die
		Subgrapheigenschaft gilt: $(u_i, u_j) \in E_N \Rightarrow (\phi(u_i), \phi(u_j)) \in E_S$.
		Diese Prüfung erfolgt durch Adjazenzmatrix-Lookup in $O(1)$ pro Kante.
	\end{enumerate}

	\begin{figure}[H]
		\centering
		\begin{tikzpicture}[
			qnode/.style={circle, draw=mainblue, fill=mainblue!20, minimum size=0.7cm,
				font=\small\bfseries},
			hnode/.style={circle, draw=darkgreen, fill=darkgreen!15, minimum size=0.7cm,
				font=\small\bfseries},
			mnode/.style={circle, draw=accentred, fill=accentred!20, minimum size=0.7cm,
				font=\small\bfseries},
			qarr/.style={-{Stealth[length=5pt]}, mainblue, thick},
			harr/.style={-{Stealth[length=5pt]}, darkgreen, thick},
		]
			% Query graph G_N (left)
			\node[font=\small\bfseries] at (1.4, 3.2) {Query $G_N$};
			\node[qnode] (u1) at (0.6, 2.3) {$u_1$};
			\node[qnode] (u2) at (2.2, 2.3) {$u_2$};
			\node[qnode] (u3) at (1.4, 1.1) {$u_3$};
			\draw[qarr] (u1) -- (u2);
			\draw[qarr] (u2) -- (u3);
			\draw[qarr] (u1) -- (u3);

			% Signatures
			\node[font=\scriptsize, darkgray] at (0.2, 1.9)
				{$\sigma\!=\!11$};
			\node[font=\scriptsize, darkgray] at (2.6, 1.9)
				{$\sigma\!=\!7$};
			\node[font=\scriptsize, darkgray] at (1.9, 0.7)
				{$\sigma\!=\!13$};

			% Host graph G_S (right)
			\node[font=\small\bfseries] at (6.0, 3.2) {Host $G_S$ (Treibergraph)};
			\node[hnode] (v1) at (4.5, 2.5) {$v_1$};
			\node[hnode] (v2) at (6.0, 3.0) {$v_2$};
			\node[hnode] (v3) at (7.5, 2.5) {$v_3$};
			\node[hnode] (v4) at (5.2, 1.1) {$v_4$};
			\node[mnode] (v5) at (6.8, 1.1) {$v_5$};
			\node[hnode] (v6) at (6.0, 0.1) {$v_6$};
			\draw[harr] (v1) -- (v2);
			\draw[harr] (v2) -- (v3);
			\draw[harr] (v1) -- (v4);
			\draw[harr] (v4) -- (v5);
			\draw[harr] (v5) -- (v3);
			\draw[harr, accentred, very thick] (v1) -- (v5);
			\draw[harr, accentred, very thick] (v5) -- (v6);
			\draw[harr, accentred, very thick] (v1) -- (v6);

			% Mapping arrows
			\draw[-{Stealth[length=4pt]}, dashed, darkgray, thick]
				(u1.east) to[bend left=12] node[above, font=\scriptsize]{$\phi(u_1)$} (v1.west);
			\draw[-{Stealth[length=4pt]}, dashed, darkgray, thick]
				(u2.south east) to[bend left=5] node[above, font=\scriptsize]{$\phi(u_2)$} (v5.north west);
			\draw[-{Stealth[length=4pt]}, dashed, darkgray, thick]
				(u3.east) to[bend right=10] node[below, font=\scriptsize]{$\phi(u_3)$} (v6.west);

			\node[font=\scriptsize, darkgray, align=center] at (6.0, -0.6)
				{Rote Kanten: gefundener Subgraph $G_N \subseteq G_S$};
		\end{tikzpicture}
		\caption{Subgraph-Isomorphismus-Beispiel für Treiber-Kompatibilitätsprüfung.
			Der Query-Graph $G_N$ (blau, links) repräsentiert die Abhängigkeiten
			eines neuen Treibers ($u_1 \to u_2$, $u_2 \to u_3$, $u_1 \to u_3$).
			Der Host-Graph $G_S$ (grün, rechts) ist der aktuelle Treibergraph
			des Systems. Die gestrichelten Pfeile zeigen die Abbildung $\phi$:
			$u_1 \mapsto v_1$, $u_2 \mapsto v_5$, $u_3 \mapsto v_6$. Die
			rot markierten Kanten in $G_S$ bilden den Subgraphen $G_N$ ab.
			Der neue Treiber ist damit kompatibel.}
		\label{fig:subgraph-ex}
	\end{figure}
	
	\begin{satz}[Subgraph Algorithmus Komplexität]\label{satz:subgraph_comp}
		Der Subgraph Algorithmus hat Gesamtkomplexität $O(n^3)$ im Worst Case
		und $O(n^2)$ im Average Case mit Signaturfilterung.
	\end{satz}
	
	\begin{proof}
		Ohne Filterung entspricht Subgraph-Isomorphismus dem NP-vollständigen
		Problem \cite{Cook1971}. Der Subgraph Algorithmus nutzt jedoch folgende
		Struktur: (1) Signaturberechnung in $O(n)$; (2) Kandidatenfilterung
		reduziert die Kandidatenmenge auf $O(\sqrt{n})$ im Average Case;
		(3) Adjazenzprüfung für jeden Kandidaten in $O(n)$. Gesamt: $O(n^2)$
		im Average Case. Im Worst Case (alle Signaturen gleich) degeneriert
		dies zu $O(n^3)$ durch vollständige Adjazenzprüfung.
	\end{proof}
	
	\section{Boolesche Matrixmultiplikation}
	
	\begin{definition}[Boolesche Matrix]
		Eine \textbf{Boolesche Matrix} $A \in \{0,1\}^{n \times n}$ hat
		Einträge $a_{ij} \in \{0,1\}$ mit Boolescher Addition
		($0 \oplus 0 = 0$, $0 \oplus 1 = 1 \oplus 0 = 1 \oplus 1 = 1$)
		und Multiplikation ($0 \otimes 0 = 0 \otimes 1 = 1 \otimes 0 = 0$,
		$1 \otimes 1 = 1$).
	\end{definition}
	
	\begin{satz}[Transitive Hülle via BMM]\label{satz:bmm}
		Die transitive Hülle $G^+$ eines Graphen $G$ mit Adjazenzmatrix $A$
		kann in $O(n^{2.37})$ berechnet werden \cite{Epp2026boolmm}.
	\end{satz}
	
	Im OS-Kontext dient die transitive Hülle zur Verifikation von
	Treiber-Abhängigkeitsgraphen: Falls $G^+_{ij} = 1$, hängt Treiber~$i$
	(transitiv) von Treiber~$j$ ab.
	
	\section{Treiber-Abhängigkeitsgraph}
	
	\begin{definition}[Treiber-Abhängigkeitsgraph]
		Der \textbf{Treiber-Abhängigkeitsgraph} $G_D = (V_D, E_D)$ hat als
		Knoten die Treiber $V_D = \{d_1, \ldots, d_m\}$ und als Kanten
		$E_D = \{(d_i, d_j) \mid d_i \text{ benötigt } d_j\}$.
	\end{definition}
	
	Satz~14.2 beweist formal, dass die Treiber-Kompatibilitätsprüfung
	auf Subgraph-Isomor\-phismus in $G_D$ reduzierbar ist und in $O(n^3)$
	entschieden werden kann.
	
	% ==============================================================
	%  KAPITEL 9 -- SYSTEMZUSTÄNDE UND LYAPUNOV-STABILITÄT
	% ==============================================================
	\chapter{Systemzustände und Lyapunov-Stabilität}
	
	\section{Systemzustands-Klassifikation}

	\begin{definition}[Systemzustand]
		Der \textbf{Systemzustand} des BS/EDF$^+$ zum Zeitpunkt $t$ ist der
		Vektor
		\[
		\mathbf{x}_t = (u_t,\, m_t,\, i_t,\, l_t)^T \in \mathbb{R}^4,
		\]
		wobei $u_t \in [0,1]$ die CPU-Auslastung, $m_t \in [0,1]$ die
		Speicherauslastung, $i_t \in [0,1]$ die I/O-Auslastung und
		$l_t \in [0,1]$ die Netzwerkauslastung bezeichnen.
	\end{definition}

	\begin{definition}[Systemzustands-Klassen KS-I bis KS-IV]
		Analog zu \cite{Epp2026sysstate} klassifizieren wir:
		\begin{itemize}
			\item \textbf{KS-I} (Nominal): $\|\mathbf{x}_t\|_\infty < 0{,}7$
			\item \textbf{KS-II} (Erhöht): $0{,}7 \leq \|\mathbf{x}_t\|_\infty < 0{,}85$
			\item \textbf{KS-III} (Kritisch): $0{,}85 \leq \|\mathbf{x}_t\|_\infty < 1{,}0$
			\item \textbf{KS-IV} (Ãœberlast): $\|\mathbf{x}_t\|_\infty = 1{,}0$
		\end{itemize}
	\end{definition}

	Abbildung~\ref{fig:ks-automat} zeigt den Zustandsautomaten der vier
	Klassen mit den Ãœbergangsbedingungen. Auf BS/EDF$^+$ reagiert das
	System auf Klassenübergänge durch adaptives Scheduling: In KS-II wird
	die DVFS-Frequenz um $10\,\%$ angehoben; in KS-III um weitere $20\,\%$.
	In KS-IV wird ein Notfall-Scheduling-Modus aktiviert, der alle
	Best-Effort-Tasks suspendiert.

	\begin{figure}[H]
		\centering
		\begin{tikzpicture}[
			state/.style={ellipse, draw=mainblue, fill=mainblue!12, minimum width=2.2cm,
				minimum height=1.0cm, align=center, font=\small\bfseries},
			crit/.style={ellipse, draw=accentred, fill=accentred!15, minimum width=2.2cm,
				minimum height=1.0cm, align=center, font=\small\bfseries},
			arr/.style={-{Stealth[length=5pt]}, thick},
			lbl/.style={font=\scriptsize, align=center, text width=2.4cm}
		]
			\node[state] (ks1) at (0, 0)    {\textbf{KS-I}\\Nominal};
			\node[state] (ks2) at (4, 0)    {\textbf{KS-II}\\Erhöht};
			\node[crit]  (ks3) at (8, 0)    {\textbf{KS-III}\\Kritisch};
			\node[crit]  (ks4) at (8, -2.8) {\textbf{KS-IV}\\Ãœberlast};

			% KS-I -> KS-II
			\draw[arr, darkgray] (ks1.east) to[bend left=18]
				node[above, lbl]{$\|\mathbf{x}\|_\infty \geq 0.7$} (ks2.west);
			% KS-II -> KS-I
			\draw[arr, darkgreen] (ks2.west) to[bend left=18]
				node[below, lbl]{$\|\mathbf{x}\|_\infty < 0.7$} (ks1.east);
			% KS-II -> KS-III
			\draw[arr, darkgray] (ks2.east) to[bend left=18]
				node[above, lbl]{$\|\mathbf{x}\|_\infty \geq 0.85$} (ks3.west);
			% KS-III -> KS-II
			\draw[arr, darkgreen] (ks3.west) to[bend left=18]
				node[below, lbl]{$\|\mathbf{x}\|_\infty < 0.85$} (ks2.east);
			% KS-III -> KS-IV
			\draw[arr, accentred] (ks3.south) to[bend left=10]
				node[right, lbl, accentred]{$\|\mathbf{x}\|_\infty = 1.0$} (ks4.north);
			% KS-IV -> KS-III
			\draw[arr, darkgreen] (ks4.north) to[bend left=10]
				node[left, lbl]{$\|\mathbf{x}\|_\infty < 1.0$} (ks3.south);

			% Self-loops
			\draw[arr, darkgray] (ks1.north west) to[loop left, looseness=5]
				node[left, font=\scriptsize]{$< 0.7$} (ks1.south west);
			\draw[arr, accentred] (ks4.east) to[loop right, looseness=5]
				node[right, font=\scriptsize]{$= 1.0$} (ks4.east);

			% DVFS annotations
			\node[font=\scriptsize, darkgray] at (4, -2.0)
				{EDF$^+$-Reaktion pro Klasse:};
			\node[font=\scriptsize] at (4, -2.5)
				{KS-I: Normal, KS-II: +10\,\%\,f, KS-III: +30\,\%\,f, KS-IV: Notfallmodus};
		\end{tikzpicture}
		\caption{Zustandsautomat der Systemzustands-Klassen KS-I bis KS-IV.
			Grüne Pfeile symbolisieren Entspannung (abnehmende Auslastung),
			graue und rote Pfeile Eskalation. BS/EDF$^+$ reagiert auf jeden
			Klassenübergang durch Anpassung der DVFS-Frequenz und ggf.\ Aktivierung
			des Notfall-Scheduling-Modus (KS-IV), in dem alle Best-Effort-Tasks
			suspendiert und nur noch EDF$^+$-Echtzeit-Tasks ausgeführt werden.}
		\label{fig:ks-automat}
	\end{figure}

	\paragraph{Klassifikations-Algorithmus.}
	Der Monitor-Task $T_{\mathrm{mon}}$ (EDF$^+$-Task, Periode $P = 10\,\text{ms}$)
	berechnet in jedem Durchlauf den aktuellen Systemzustand~$\mathbf{x}_t$:
	CPU-Auslastung aus dem Kernel-Scheduler-Statistikpuffer, Speicherauslastung
	aus dem Buddy-Allokator, I/O-Auslastung aus der Block-Device-Layer-Queue,
	Netzwerkauslastung aus der Treiber-TX/RX-Statistik. Die Infinity-Norm
	$\|\mathbf{x}_t\|_\infty$ wird bestimmt und die Zustandsklasse nachgeführt.
	Klassenübergänge lösen asynchrone Callbacks in der EDF$^+$-Scheduler-Core
	aus, die die DVFS-Parameter des IoFET-Controllers aktualisieren.

	\section{Dynamisches System und Zustandsmatrix}
	
	\begin{definition}[OS-Zustandsübergangsmatrix]
		Das OS-Systemzustand-Dynamikmodell ist
		\[
		\mathbf{x}_{t+1} = A \mathbf{x}_t + B \mathbf{u}_t,
		\]
		wobei $A \in \mathbb{R}^{4 \times 4}$ die \textbf{Systemmatrix},
		$B \in \mathbb{R}^{4 \times k}$ die Eingangsmatrix und
		$\mathbf{u}_t$ der EDF$^+$-Scheduling-Eingangsvektor ist.
	\end{definition}
	
	\begin{definition}[Lyapunov-Funktion]
		Eine Funktion $V: \mathbb{R}^n \to \mathbb{R}_{\geq 0}$ heißt
		\textbf{Lyapunov-Funktion} für $\mathbf{x}_{t+1} = A\mathbf{x}_t$, falls:
		\begin{enumerate}
			\item $V(\mathbf{0}) = 0$ und $V(\mathbf{x}) > 0$ für $\mathbf{x} \neq \mathbf{0}$,
			\item $V(A\mathbf{x}) < V(\mathbf{x})$ für alle $\mathbf{x} \neq \mathbf{0}$.
		\end{enumerate}
	\end{definition}
	
	\begin{satz}[Stabilität durch Lyapunov]\label{satz:lyap_prelim}
		Das System $\mathbf{x}_{t+1} = A\mathbf{x}_t$ ist asymptotisch stabil
		genau dann, wenn der Spektralradius
		$\rho(A) = \max_i |\lambda_i(A)| < 1$.
	\end{satz}
	
	\begin{proof}
		Die Jordanzerlegung $A = P J P^{-1}$ liefert $A^t = P J^t P^{-1}$.
		Falls $\rho(A) < 1$, gilt $\|J^t\| \to 0$ für $t \to \infty$.
		Daher $A^t \to 0$, also $\mathbf{x}_t \to \mathbf{0}$.
		Umgekehrt: falls ein $|\lambda_i| \geq 1$, wächst der entsprechende
		Eigenvektor-Anteil nicht gegen Null.
	\end{proof}
	
	\section{Steuerbarkeit und Beobachtbarkeit}
	
	\begin{definition}[Steuerbarkeit]
		Das System $(A, B)$ ist \textbf{steuerbar}, falls
		$\rank[B \mid AB \mid \cdots \mid A^{n-1}B] = n$.
	\end{definition}
	
	\begin{definition}[Beobachtbarkeit]
		Das System $(A, C)$ ist \textbf{beobachtbar}, falls
		$\rank[C^T \mid A^T C^T \mid \cdots \mid (A^T)^{n-1}C^T]^T = n$.
	\end{definition}
	
	% ==============================================================
	%  KAPITEL 10 -- DATEISYSTEM-INITIALISIERUNG
	% ==============================================================
	\chapter{Dateisystem-Initialisierung}
	
	\section{Inode-Graph und Deduplizierung}

	\begin{definition}[Inode-Graph]
		Der \textbf{Inode-Graph} $G_{\mathrm{FS}} = (V_{\mathrm{FS}}, E_{\mathrm{FS}})$
		hat als Knoten die Inodes (Dateien, Verzeichnisse) und als Kanten
		die Verzeichnis-Einträge.
	\end{definition}

	Abbildung~\ref{fig:inode-graph} zeigt einen beispielhaften Inode-Graphen
	mit drei Verzeichnissen und vier Dateien, von denen zwei identische
	Inhalte besitzen (Inodes~$F_3$ und $F_4$). Das Deduplizierungssystem
	erkennt diese Gleichheit über SHA3-256-Hashing und ersetzt eine der
	Kopien durch einen Hard-Link.

	\begin{figure}[H]
		\centering
		\begin{tikzpicture}[
			dir/.style={rectangle, draw=mainblue, fill=mainblue!15, rounded corners=3pt,
				minimum width=1.5cm, minimum height=0.7cm, align=center, font=\small},
			file/.style={rectangle, draw=darkgreen, fill=darkgreen!12, rounded corners=3pt,
				minimum width=1.7cm, minimum height=0.7cm, align=center, font=\small},
			dupfile/.style={rectangle, draw=accentred, fill=accentred!15, rounded corners=3pt,
				minimum width=1.7cm, minimum height=0.7cm, align=center, font=\small},
			arr/.style={-{Stealth[length=4pt]}, darkgray, thick},
			hl/.style={-{Stealth[length=4pt]}, accentred, thick, dashed}
		]
			% Root inode
			\node[dir] (root) at (4.5, 4.5) {/ (Root)};

			% Subdirectories
			\node[dir] (home) at (1.5, 3.0) {/home};
			\node[dir] (etc)  at (4.5, 3.0) {/etc};
			\node[dir] (tmp)  at (7.5, 3.0) {/tmp};

			% Files
			\node[file]    (f1) at (0.3, 1.4) {$F_1$\\config.txt};
			\node[file]    (f2) at (2.7, 1.4) {$F_2$\\data.bin};
			\node[dupfile] (f3) at (4.5, 1.4) {$F_3$\\cache.dat};
			\node[dupfile] (f4) at (7.5, 1.4) {$F_4$\\cache.dat};

			% Edges
			\draw[arr] (root) -- (home);
			\draw[arr] (root) -- (etc);
			\draw[arr] (root) -- (tmp);
			\draw[arr] (home) -- (f1);
			\draw[arr] (home) -- (f2);
			\draw[arr] (etc)  -- (f3);
			\draw[arr] (tmp)  -- (f4);

			% Hard link after dedup
			\draw[hl] (f4.south) to[bend right=-30]
				node[below, align=center, font=\scriptsize, text=accentred]{Hard-Link\\(nach Dedup)} (f3.south);

			% SHA3 labels
			\node[font=\scriptsize, accentred] at (4.5, 0.5)
				{SHA3\,=\,\texttt{0xAB3F...}};
			\node[font=\scriptsize, accentred] at (7.5, 0.5)
				{SHA3\,=\,\texttt{0xAB3F...}};
		\end{tikzpicture}
		\caption{Inode-Graph eines BS/EDF$^+$-Dateisystems mit drei Verzeichnissen
			und vier Dateien. Die Dateien $F_3$ (\texttt{/etc/cache.dat}) und
			$F_4$ (\texttt{/tmp/cache.dat}) haben identische SHA3-256-Hashes
			(rot markiert). Der Deduplizierungsalgorithmus ersetzt $F_4$ durch
			einen Hard-Link auf $F_3$ (roter gestrichelter Pfeil), sodass beide
			Verzeichniseinträge auf denselben Inode zeigen. Im Inode-Graphen
			entsteht damit eine Kante von \texttt{/tmp} direkt zu $F_3$.}
		\label{fig:inode-graph}
	\end{figure}

	Das Deduplizierungssystem \cite{Epp2026dusgraph} identifiziert
	isomorphe Teilgraphen im $G_{\mathrm{FS}}$. Der Subgraph Algorithmus
	\cite{Epp2026subgraph} liefert das algorithmische Rückgrat.

	\paragraph{Deduplizierungs-Algorithmus.}
	Algorithmus~\ref{alg:dedup} führt die Dateisystem-Deduplizierung durch.
	In Phase~1 (Zeilen~2--4) werden für alle Inodes SHA3-256-Hashes berechnet.
	SHA3-256 ist kollisionsresistent mit einer Kollisionswahrscheinlichkeit
	von $2^{-128}$ und daher für praktische Zwecke als injektiv anzusehen.
	In Phase~2 (Zeilen~5--10) werden Inodes nach gleichem Hash gruppiert.
	Jede Gruppe mit mehr als einem Element enthält duplizierten Inhalt.
	Für jede solche Gruppe wird ein Repräsentant~$v_{\mathrm{repr}}$ gewählt
	(derjenige Inode mit dem ältesten Erstellungszeitpunkt), und alle anderen
	Inodes der Gruppe werden durch Hard-Links auf~$v_{\mathrm{repr}}$ ersetzt.
	Hard-Links sind OS-seitig $O(1)$: der Verzeichniseintrag wird einfach
	auf die Inode-Nummer von~$v_{\mathrm{repr}}$ geändert.

	\begin{algorithm}[H]
		\caption{Dateisystem-Deduplizierung}
		\label{alg:dedup}
		\begin{algorithmic}[1]
			\Procedure{FsDeduplicate}{$G_{\mathrm{FS}}$}
			\For{each Inode $v \in V_{\mathrm{FS}}$}
			\State $h(v) \gets \mathrm{SHA3\_256}(\mathrm{content}(v))$
			\EndFor
			\State $\mathrm{Groups} \gets \mathrm{GroupBy}(V_{\mathrm{FS}}, h)$
			\For{each group $G_k$ with $|G_k| > 1$}
			\State $v_{\mathrm{repr}} \gets G_k[0]$
			\For{each $v \in G_k \setminus \{v_{\mathrm{repr}}\}$}
			\State \Call{HardLink}{$v$, $v_{\mathrm{repr}}$}
			\EndFor
			\EndFor
			\EndProcedure
		\end{algorithmic}
	\end{algorithm}
	
	\section{Mobile-OS-Dateiarchitektur}
	
	Für mobile Geräte nutzt BS/EDF$^+$ eine partitionierte Dateisystem-Architektur
	\cite{Epp2026mobde}:
	\begin{itemize}
		\item \textbf{/system}: Read-Only, A/B-OTA-Update-fähig.
		\item \textbf{/data}: Verschlüsselt (AES-256-XTS), benutzerspezifisch.
		\item \textbf{/cache}: Temporär, automatisch bereinigt durch
		$T_{\mathrm{cache\_evict}}$ (EDF$^+$-Task, nie\-drige Priorität).
	\end{itemize}
	
	% ==============================================================
	%  KAPITEL 11 -- GERÄTETREIBER-ARCHITEKTUR
	% ==============================================================
	\chapter{Gerätetreiber-Architektur}
	
	\section{AUTOSAR-Treiber-Framework}

	AUTOSAR \cite{AUTOSAR2023} definiert eine mehrschichtige Softwarearchitektur:
	\begin{itemize}
		\item \textbf{MCAL}: Direkter Hardware-Zugriff (SPI, CAN, LIN, ETH, PWM).
		\item \textbf{ECU-AL}: HAL über MCAL.
		\item \textbf{Services Layer}: OS, Memory, Communication.
		\item \textbf{RTE}: Verbindet SWCs.
		\item \textbf{SWC}: Anwendungslogik.
	\end{itemize}

	Abbildung~\ref{fig:autosar-layers} zeigt den AUTOSAR-Schichtenaufbau und
	seine Integration in BS/EDF$^+$. Jede Schicht kommuniziert über klar
	definierte Interfaces mit den angrenzenden Schichten; direkte Schicht-Sprünge
	sind unzulässig.

	\begin{figure}[H]
		\centering
		\begin{tikzpicture}[
			layer/.style={rectangle, draw=mainblue, fill=mainblue!10,
				minimum width=9cm, minimum height=0.75cm, align=center, font=\small},
			hlayer/.style={rectangle, draw=accentred, fill=accentred!12,
				minimum width=9cm, minimum height=0.75cm, align=center, font=\small},
			hwlayer/.style={rectangle, draw=darkgray, fill=darkgray!15,
				minimum width=9cm, minimum height=0.75cm, align=center, font=\small},
			darr/.style={{Stealth[length=4pt]}-{Stealth[length=4pt]}, thick, darkgray}
		]
			\node[layer]   at (0, 6.4) {\textbf{SWC} -- Anwendungssoftware-Komponenten};
			\node[layer]   at (0, 5.5) {RTE (Runtime Environment)};
			\node[hlayer]  at (0, 4.6) {\textbf{Services Layer:} OS (BS/EDF$^+$), COM, MEM};
			\node[layer]   at (0, 3.7) {\textbf{ECU Abstraction Layer (ECU-AL)}};
			\node[hlayer]  at (0, 2.8) {\textbf{MCAL:} SPI, CAN, LIN, ADC, PWM, ETH};
			\node[hwlayer] at (0, 1.9) {Mikrocontroller / ECU-Hardware};

			% Interface arrows
			\draw[darr] (4.6, 6.05) -- (4.6, 5.875);
			\draw[darr] (4.6, 5.15) -- (4.6, 4.975);
			\draw[darr] (4.6, 4.25) -- (4.6, 4.075);
			\draw[darr] (4.6, 3.35) -- (4.6, 3.175);
			\draw[darr] (4.6, 2.45) -- (4.6, 2.275);

			% EDF+ annotation
			%\draw[-{Stealth}, accentred, thick] (4.9, 4.6) -- (5.9, 4.6)
			%	node[right, font=\scriptsize, accentred]{EDF$^+$-Scheduler};
		\end{tikzpicture}
		\caption{AUTOSAR-Schichtenarchitektur und BS/EDF$^+$-Integration.
			Die Anwendungssoftware-Komponenten (SWC, oberste Schicht) kommunizieren
			nur über die Laufzeitumgebung (RTE) mit der restlichen Architektur.
			Der BS/EDF$^+$-Kernel ist im Services Layer eingebettet (rot markiert)
			und steuert alle Tasks über EDF$^+$. Die MCAL-Schicht kapselt den
			direkten Hardware-Zugriff; Treiber-Abhängigkeiten werden via
			Subgraph Algorithmus (Kapitel~8) auf Kompatibilität geprüft, bevor
			ein MCAL-Modul geladen wird.}
		\label{fig:autosar-layers}
	\end{figure}

	\begin{definition}[ECU-Zustandsmaschine]
		Die \textbf{ECU-Zustandsmaschine} hat Zustände
		$Q = \{Q_{\mathrm{init}}, Q_{\mathrm{run}}, Q_{\mathrm{sleep}},
		Q_{\mathrm{fault}}, Q_{\mathrm{shutdown}}\}$ und Ãœbergangsfunktion
		$\delta: Q \times \Sigma \to Q$.
	\end{definition}

	BS/EDF$^+$ implementiert AUTOSAR-OTA-Updates \cite{Epp2026autsre}
	als EDF$^+$-Task $T_{\mathrm{OTA}}$ mit Deadline $D_{\mathrm{OTA}}$
	nach ISO~26262 ASIL. Der OTA-Update-Prozess verläuft wie folgt:
	Zunächst wird das Update-Image über das Netzwerk (EtherCAT oder ETH)
	empfangen und in einer Schattierpartition (\emph{B-Slot}) gespeichert.
	Dann prüft ein Integritäts-Task $T_{\mathrm{verify}}$ die Signatur
	des Images mit CRYSTALS-Dilithium-3. Besteht die Prüfung, wird der
	Bootmanager angewiesen, beim nächsten Start aus dem B-Slot zu booten.
	Schlägt die Verifikation fehl, verbleibt das System im A-Slot (aktuelle
	Produktionsversion) und eine Fehlermeldung wird protokolliert.

	\section{EtherCAT-Integration}

	EtherCAT \cite{IEC61158,Epp2026ethercate} ist echtzeitfähiges
	Ethernet mit $\leq 100\,\mu\text{s}$ Zykluszeit.

	\begin{definition}[EtherCAT-Telegramm]
		$F_{\mathrm{EC}} = (\mathrm{Header},\, \mathrm{Cmd},\, \mathrm{Addr},\,
		\mathrm{Data}[1..n],\, \mathrm{WKC})$ mit Working Counter WKC.
	\end{definition}

	\paragraph{EtherCAT-Frame-Verarbeitungs-Algorithmus.}
	Im EHAE-Framework \cite{Epp2026ethercate} wird jedes eingehende EtherCAT-Telegramm
	als EDF$^+$-Task mit harter Deadline $D_{\mathrm{EC}} = T_{\mathrm{Zyklus}}$
	registriert. Der Verarbeitungsalgorithmus arbeitet wie folgt: Beim Eingang
	eines Ethernet-Frames prüft der NIC-Interrupt-Handler zunächst den
	EtherType-Wert (\texttt{0x88A4} für EtherCAT). Wenn erkannt, wird ein
	neuer EDF$^+$-Task $T_{\mathrm{EC}}$ erstellt mit Deadline
	$D = t_{\mathrm{now}} + T_{\mathrm{Zyklus}}$. Dieser Task durchläuft das
	Telegramm sequenziell: Jeder EtherCAT-Slave liest seinen Datenbereich,
	schreibt seine Prozessdaten, inkrementiert den Working Counter und leitet
	das Telegramm weiter. Nach dem letzten Slave wird das Telegramm an den
	Master zurückgesendet. Der Master-Task prüft, ob $\mathrm{WKC}$ dem
	erwarteten Wert entspricht (Anzahl aktiver Slaves); Abweichungen signalisieren
	einen ausgefallenen Slave.
	
	\section{Python-Treiber-Framework}
	
	\begin{definition}[Treiber-Interface]
		Ein Python-Treiber implementiert:
		\begin{lstlisting}[language=Python]
class Driver:
	def probe(self, device): ...
	def init(self) -> None: ...
	def read(self, reg: int) -> int: ...
	def write(self, reg: int, val: int) -> None: ...
	def cleanup(self) -> None: ...
		\end{lstlisting}
	\end{definition}
	
	% ==============================================================
	%  KAPITEL 12 -- ENERGIEMANAGEMENT
	% ==============================================================
	\chapter{Energiemanagement}
	
	\section{IoFET-Sigmoid-DVFS}

	\begin{definition}[DVFS-Steuergesetz]
		Das \textbf{Dynamic Voltage and Frequency Scaling} folgt dem
		IoFET-Modell \cite{Epp2026iono}:
		\[
		V^*(f) = V_{\min} + (V_{\max} - V_{\min}) \cdot
		\sigma\!\left(\frac{f - f_{1/2}}{\Delta f}\right),
		\quad
		\sigma(x) = \frac{1}{1 + e^{-x}}.
		\]
	\end{definition}

	Die IoFET-Sigmoid-Kurve beschreibt, wie die Betriebsspannung mit
	steigender Frequenz sanft und monoton ansteigt. Abbildung~\ref{fig:dvfs-sigmoid}
	zeigt die Kurve für typische Parameterwerte ($V_{\min} = 0.7\,\text{V}$,
	$V_{\max} = 1.3\,\text{V}$, $f_{1/2} = 3\,\text{GHz}$, $\Delta f = 1\,\text{GHz}$).
	Im Betrieb schlägt der IoFET-DVFS-Controller nach jedem EDF$^+$-Scheduling-Zyklus
	die optimale Betriebsspannung nach, indem er die aktuelle Taktfrequenz $f(t)$
	in die Sigmoid-Formel einsetzt und das Ergebnis in das MSR~\texttt{IA32\_PERF\_CTL}
	schreibt.

	\begin{figure}[H]
		\centering
		\begin{tikzpicture}
			\begin{axis}[
				width=10cm, height=6cm,
				xlabel={Taktfrequenz $f$ [GHz]},
				ylabel={Betriebsspannung $V^*(f)$ [V]},
				xmin=0, xmax=6, ymin=0.55, ymax=1.45,
				grid=major, grid style={lightgray!60},
				xtick={0,1,2,3,4,5,6},
				ytick={0.6,0.7,0.8,0.9,1.0,1.1,1.2,1.3,1.4},
				legend pos=north west,
			]
				% Sigmoid DVFS curve  V_min=0.7, V_max=1.3, f_half=3, Df=1
				\addplot[mainblue, very thick, domain=0:6, samples=300]
					{0.7 + 0.6 / (1 + exp(-(x - 3.0) / 1.0))};
				\addlegendentry{IoFET-Sigmoid $V^*(f)$}

				% Linear reference (classic DVFS)
				\addplot[darkgray, thick, dashed, domain=0:6]
					{0.7 + (1.3 - 0.7) * x / 6};
				\addlegendentry{Lineares DVFS (Referenz)}

				% Optimal operating point annotation
				\addplot[accentred, mark=*, mark size=3pt] coordinates {(3.0, 1.0)};
				\node[font=\scriptsize, accentred] at (axis cs: 3.8, 0.93)
					{$(f_{1/2},\, 1.0\,\text{V})$};

				% Energy savings area (hatch)
				\addplot[fill=mainblue!15, fill opacity=0.5, draw=none,
					domain=2:5, samples=100]
					{0.7 + 0.6 / (1 + exp(-(x-3)/1))} \closedcycle;
			\end{axis}
		\end{tikzpicture}
		\caption{IoFET-Sigmoid-DVFS-Kurve (blau) im Vergleich zu einem linearen
			DVFS-Referenzmodell (gestrichelt). Die Sigmoid-Kurve hält die
			Betriebsspannung bei niedrigen Frequenzen nahe $V_{\min} = 0.7\,\text{V}$
			und steigt erst ab dem Wendepunkt $f_{1/2} = 3\,\text{GHz}$ steil an.
			Gegenüber dem linearen Modell spart die Sigmoid-Kurve im Niedriglast-Bereich
			($f < 2\,\text{GHz}$) bis zu $18\,\%$ Verlustleistung (blau schraffierter
			Bereich). Korollar~\ref{kor:dvfs_opt} beweist die Energieoptimalität dieses
			Steuergesetzes.}
		\label{fig:dvfs-sigmoid}
	\end{figure}

	\begin{satz}[DVFS-Konvexität]\label{satz:dvfs_convex}
		Die Energiefunktion $E(f) = \alpha C_L V^*(f)^2$ ist konvex in $f$
		für das IoFET-Steuergesetz.
	\end{satz}

	\begin{proof}
		$V^*(f)$ ist eine glatte, streng monoton wachsende Funktion von $f$
		(Sigmoid ist konvex für $f > f_{1/2}$). Das Quadrat einer konvexen
		nicht-negativen Funktion ist wieder konvex. Also ist $E(f)$ konvex,
		und das eindeutige Minimum wird durch $dE/df = 0$ bestimmt.
	\end{proof}

	\section{Li-Ionen-Akkumodell}

	\begin{definition}[Äquivalentkreis-Modell]
		Die Klemmenspannung des Thevenin-Ersatz\-schaltbilds \cite{Epp2026liionp}:
		\[
		V(t) = V_{\mathrm{OCV}}(\mathrm{SoC}) - I(t) R_0 - V_{R_1 C_1}(t) - V_{R_2 C_2}(t).
		\]
	\end{definition}

	Abbildung~\ref{fig:liion-circuit} zeigt das Thevenin-Ersatzschaltbild
	des Li-Ionen-Akku-Modells. Die offene Klemmspannung $V_{\mathrm{OCV}}$ hängt
	nichtlinear vom Ladezustand (State of Charge, SoC) ab und wird in BS/EDF$^+$
	durch eine Lookup-Table mit 256 Stützpunkten approximiert.

	\begin{figure}[H]
		\centering
		\begin{tikzpicture}[
			circuit/.style={draw=darkgray, thick},
			lbl/.style={font=\small}
		]
			% Battery OCV source
			\draw[circuit] (0,0) -- (0,1.0);
			\draw[circuit] (-0.25,1.0) -- (0.25,1.0);   % minus
			\draw[circuit] (-0.4,1.3) -- (0.4,1.3);    % plus (thicker)
			\draw[circuit, line width=2pt] (-0.4,1.3) -- (0.4,1.3);
			\draw[circuit] (0,1.3) -- (0,1.8);
			\node[lbl, left] at (-0.5, 1.15) {$V_{\mathrm{OCV}}$};

			% R0 (series resistance)
			\draw[circuit] (0,1.8) -- (1.2,1.8);
			\draw[circuit] (1.2,1.4) rectangle (1.8,2.2);
			\node[lbl, above] at (1.5,2.2) {$R_0$};
			\draw[circuit] (1.8,1.8) -- (3.0,1.8);

			% R1C1 parallel
			\draw[circuit] (3.0,1.8) -- (3.0,2.6);
			\draw[circuit] (3.0,2.6) -- (3.6,2.6);
			\draw[circuit] (3.6,2.2) rectangle (4.2,3.0);
			\node[lbl, above] at (3.9,3.0) {$R_1$};
			\draw[circuit] (4.2,2.6) -- (4.8,2.6);

			\draw[circuit] (3.0,1.8) -- (3.0,1.0);
			\draw[circuit] (3.0,1.0) -- (3.7,1.0);
			% Capacitor C1
			\draw[circuit] (3.7,0.7) -- (3.7,1.3);
			\draw[circuit] (4.1,0.7) -- (4.1,1.3);
			\node[lbl, below] at (3.9,0.7) {$C_1$};
			\draw[circuit] (4.1,1.0) -- (4.8,1.0);
			\draw[circuit] (4.8,1.0) -- (4.8,2.6);

			% R2C2 parallel
			\draw[circuit] (4.8,1.8) -- (5.5,1.8);
			\draw[circuit] (5.5,1.8) -- (5.5,2.6);
			\draw[circuit] (5.5,2.6) -- (6.1,2.6);
			\draw[circuit] (6.1,2.2) rectangle (6.7,3.0);
			\node[lbl, above] at (6.4,3.0) {$R_2$};
			\draw[circuit] (6.7,2.6) -- (7.3,2.6);

			\draw[circuit] (5.5,1.8) -- (5.5,1.0);
			\draw[circuit] (5.5,1.0) -- (6.2,1.0);
			\draw[circuit] (6.2,0.7) -- (6.2,1.3);
			\draw[circuit] (6.6,0.7) -- (6.6,1.3);
			\node[lbl, below] at (6.4,0.7) {$C_2$};
			\draw[circuit] (6.6,1.0) -- (7.3,1.0);
			\draw[circuit] (7.3,1.0) -- (7.3,2.6);

			% Terminal
			\draw[circuit] (7.3,1.8) -- (8.2,1.8);
			\draw[circuit] (0,0) -- (8.2,0);
			\draw[circuit] (8.2,0) -- (8.2,1.8);
			\node[lbl, right] at (8.2,0.9) {$V_T$};
			\fill[darkgray] (8.2,1.8) circle (2pt);
			\fill[darkgray] (8.2,0) circle (2pt);

			% Current arrow
			\draw[-{Stealth}, accentred, thick] (2.0,2.1) -- (2.8,2.1)
				node[above, font=\scriptsize, accentred]{$I(t)$};
		\end{tikzpicture}
		\caption{Thevenin-Ersatzschaltbild des Li-Ionen-Akku-Modells. Die
			Spannungsquelle~$V_{\mathrm{OCV}}$ modelliert die leerlaufspannungs-
			abhängige Zellspannung in Abhängigkeit vom SoC. Der Innenwiderstand~$R_0$
			repräsentiert den ohmschen Widerstand des Elektrolyten und der Elektroden.
			Die $RC$-Glieder $R_1C_1$ und $R_2C_2$ modellieren den elektrochemischen
			Doppelschicht-Effekt (langsame Diffusionsdynamik). BS/EDF$^+$ berechnet
			aus diesem Modell prädiktiv den SoC und passt die DVFS-Parameter proaktiv
			an, um unnötige Tiefentladungen zu vermeiden.}
		\label{fig:liion-circuit}
	\end{figure}

	Die SoC-Dynamik folgt:
	\[
	\frac{d\,\mathrm{SoC}}{dt} = -\frac{I(t)}{Q_{\mathrm{nom}}}.
	\]
	Das KI-gestützte PMM \cite{Epp2026liionp} prädiziert $\mathrm{SoC}(t + \Delta t)$
	und passt DVFS-Parameter proaktiv an.
	
	\section{Kontinuierliche P-Zustände}
	
	BS/EDF$^+$ implementiert kontinuierliche P-Zustände über MSR \texttt{IA32\_PERF\_CTL}:
	\[
	f(t) = f_{\min} + (f_{\max} - f_{\min}) \cdot \frac{U_{\mathrm{EDF}^+}(t)}{U_{\max}}.
	\]
	
	% ==============================================================
	%  KAPITEL 13 -- POST-QUANTEN-SICHERHEIT
	% ==============================================================
	\chapter{Post-Quanten-Sicherheit}
	
	\section{Gitterkryptographie}

	\begin{definition}[Gitter]
		Ein \textbf{Gitter} $\mathcal{L} \subset \mathbb{R}^n$ ist
		\[
		\mathcal{L} = \left\{\sum_{i=1}^m z_i \mathbf{b}_i \;\middle|\; z_i \in \mathbb{Z}\right\}.
		\]
	\end{definition}

	Abbildung~\ref{fig:lattice-lwe} zeigt ein zweidimensionales Gitter mit
	Basispunkten und verrauschten LWE-Samples. Das Grundprinzip von LWE
	besteht darin, dass die Samples $(\mathbf{a}_i, b_i)$ auf den ersten
	Blick wie zufällig aussehen -- selbst für einen Angreifer, der den
	Algorithmus kennt, ist es rechnerisch schwer, den geheimen Vektor
	$\mathbf{s}$ aus polynomiell vielen Samples $(\mathbf{a}_i, b_i)$ zu
	rekonstruieren.

	\begin{figure}[H]
		\centering
		\begin{tikzpicture}
			\begin{axis}[
				width=9cm, height=8cm,
				xlabel={$x_1$}, ylabel={$x_2$},
				xmin=-0.5, xmax=5.5, ymin=-0.5, ymax=5.5,
				grid=major, grid style={lightgray!40},
				xtick={0,1,2,3,4,5}, ytick={0,1,2,3,4,5},
				tick label style={font=\scriptsize},
			]
				% Lattice points
				\addplot[mainblue, mark=*, mark size=3.5pt,
					only marks, mark options={fill=mainblue}]
				coordinates {
					(0,0)(1,0)(2,0)(3,0)(4,0)(5,0)
					(0,1)(1,1)(2,1)(3,1)(4,1)(5,1)
					(0,2)(1,2)(2,2)(3,2)(4,2)(5,2)
					(0,3)(1,3)(2,3)(3,3)(4,3)(5,3)
					(0,4)(1,4)(2,4)(3,4)(4,4)(5,4)
					(0,5)(1,5)(2,5)(3,5)(4,5)(5,5)
				};
				\addlegendentry{Gitterpunkte $\mathcal{L}$}

				% LWE samples (noisy)
				\addplot[accentred, mark=x, mark size=5pt, only marks,
					mark options={line width=1.5pt}]
				coordinates {
					(1.15, 2.93)(2.87, 0.12)(0.08, 3.85)
					(3.92, 1.94)(4.11, 3.02)(2.03, 4.88)
					(1.92, 1.07)(3.05, 3.93)(4.88, 2.11)
				};
				\addlegendentry{LWE-Samples $(\mathbf{a}_i, b_i)$ mit Rauschen $e_i$}

				% Secret vector s
				\addplot[-{Stealth}, darkgreen, very thick] coordinates {(0,0)(2,3)};
				\node[darkgreen, font=\small] at (axis cs: 1.2, 2.0)
					{$\mathbf{s}$};
			\end{axis}
		\end{tikzpicture}
		\caption{Zwei\-dimensionales Gitter $\mathcal{L}$ (blaue Punkte) mit
			LWE-Samples (rote Kreuze). Die roten Samples liegen nahe, aber nicht
			exakt auf Gitterpunkten -- der Fehlerterm~$e_i$ verursacht eine
			kleine Verschiebung. Der geheime Vektor~$\mathbf{s}$ (grüner Pfeil)
			ist der gesuchte Parameter, den ein Angreifer aus den Samples ermitteln
			müsste. In hohen Dimensionen ($n \geq 1024$) ist dies rechnerisch
			unlösbar (SIVP-Härte nach Regev~\cite{Regev2009}).}
		\label{fig:lattice-lwe}
	\end{figure}

	\begin{definition}[LWE-Problem]
		Das \textbf{Learning With Errors}-Problem \cite{Regev2009}:
		Gegeben Samples $(\mathbf{a}_i, b_i)$ mit
		$b_i = \langle \mathbf{a}_i, \mathbf{s} \rangle + e_i \pmod{q}$
		-- finde $\mathbf{s} \in \mathbb{Z}_q^n$.
	\end{definition}

	\paragraph{LWE-Verschlüsselungs-Algorithmus.}
	Das LWE-basierte Verschlüsselungsschema für physikalische Speicheradressen
	(Satz~\ref{satz:pqc_mem}) arbeitet in drei Phasen:
	\begin{enumerate}
		\item \textbf{Schlüsselgenerierung:} Wähle zufällig
		$\mathbf{s} \gets \mathbb{Z}_q^n$ (geheimer Schlüssel) und
		generiere eine $m \times n$-Matrix $A$ gleichverteilt über $\mathbb{Z}_q^n$.
		Berechne $\mathbf{b} = A\mathbf{s} + \mathbf{e} \pmod{q}$ mit kleinem
		Fehlervektor $\mathbf{e} \in \{-\lfloor q/8 \rfloor, \ldots, \lfloor q/8 \rfloor\}^m$.
		Der öffentliche Schlüssel ist $(A, \mathbf{b})$.
		\item \textbf{Verschlüsselung einer Adresse $a$:} Wähle zufällig
		$\mathbf{r} \in \{0,1\}^m$. Berechne:
		$\mathbf{c}_1 = A^T \mathbf{r} \pmod{q}$ und
		$c_2 = \mathbf{b}^T \mathbf{r} + \lfloor q/2 \rfloor \cdot a_{\text{bit}} \pmod{q}$.
		Das Chiffrat ist $(\mathbf{c}_1, c_2)$.
		\item \textbf{Entschlüsselung:} Berechne
		$\hat{a} = c_2 - \mathbf{s}^T \mathbf{c}_1 \pmod{q}$.
		Falls $\hat{a}$ näher an $\lfloor q/2 \rfloor$ liegt als an~$0$,
		ist das Bit~$1$, sonst~$0$.
	\end{enumerate}
	Die Korrektheit folgt aus der Fehler-Kompensation: der Fehlerterm
	$\mathbf{e}^T \mathbf{r}$ bleibt klein und stört die Entschlüsselung nicht.

	\section{CRYSTALS-Kyber und Dilithium}
	
	CRYSTALS-Kyber \cite{Bos2018} ist ein IND-CCA2-sicheres KEM auf Module-LWE.
	CRYSTALS-Dilithium \cite{Ducas2018} ist ein Signaturverfahren auf Modul-LWE.
	Beide sind NIST-PQC-Standards (2024) und werden in BS/EDF$^+$ eingesetzt für:
	\begin{itemize}
		\item \textbf{Secure Boot:} Kernel-Signierung mit Dilithium-3.
		\item \textbf{TLS 1.3:} Kyber-768 als Key-Exchange.
		\item \textbf{Disk-Encryption:} Kyber-Wrapped AES-256.
	\end{itemize}
	
	\section{Quantencomputer-Bedrohungsmodell}
	
	\begin{definition}[Quantenangreifer]
		Ein \textbf{Quantenangreifer} $\mathcal{A}_Q$ kann Shors Algorithmus
		\cite{Shor1994} ausführen: RSA/ECC bricht in $O(n^3)$ Quantenoperationen.
	\end{definition}
	
	LWE mit Dimension $n \geq 1024$ liefert $128$-Bit Post-Quanten-Sicherheit.
	Selbst Grovers Algorithmus reduziert die Sicherheit nur auf $n/2$ Bits.
	
	% ==============================================================
	%  KAPITEL 14 -- NEUE ERKENNTNISSE UND FORMALE BEWEISE
	% ==============================================================
	\chapter{Neue Erkenntnisse und formale Beweise}
	
	\section{Satz 1: OS-Initialisierungssequenz als \texorpdfstring{EDF$^+$}{EDF+}-Task-System}
	
	\begin{satz}[OS-Init als EDF$^+$-Task-System]\label{satz:os_init}
		Jede Hardware-Initialisierungssequenz eines Computersystems kann als
		schedulierbares EDF$^+$-Task-System $\Sigma_{\mathrm{boot}}$ modelliert werden,
		für das EDF$^+$ die optimale Boot-Reihenfolge bestimmt.
	\end{satz}
	
	\begin{proof}
		\textbf{Konstruktion:} Definiere für jede Hardware-Komponente $h_i$
		einen Task $T_i = (C_i, D_i, P_i)$, wobei $C_i$ die Initialisierungszeit,
		$D_i$ die Deadline (durch Komponentenabhängigkeiten bestimmt) und
		$P_i = D_i$ (einmalige Ausführung).
		
		\textbf{Schedulierbarkeit:} Da jede Komponente einmal initialisiert
		wird und $C_i \leq D_i$ per Hardware-Spezifikation gilt:
		\[
		U_{\mathrm{boot}} = \sum_{i=1}^{n} \frac{C_i}{D_i} \leq 1.
		\]
		
		\textbf{Optimalität:} Nach Satz~\ref{satz:edfopt} ist EDF$^+$ optimal.
		Hardware-Abhängigkeiten $h_i \to h_j$ werden als $D_i < D_j$ kodiert,
		wodurch EDF$^+$ automatisch die korrekte Reihenfolge erzwingt.
		
		\textbf{Blocking-Behandlung:} Falls $h_i$ blockiert (Timeout), setzt
		EDF$^+$ $D_i^{\mathrm{eff}} = +\infty$ und initialisiert unabhängige
		Komponenten zuerst. Das System bleibt schedulierbar.
	\end{proof}
	
	\begin{korollar}[Boot-Parallelismus]\label{kor:boot_par}
		Auf PYMCA-8 mit $k = 8$ Kernen kann $\Sigma_{\mathrm{boot}}$ in
		$O(U_{\mathrm{boot}} \cdot t_{\max} / k)$ abgearbeitet werden.
	\end{korollar}
	
	\begin{proof}
		Partitionierung von $\Sigma_{\mathrm{boot}}$ in $k$ unabhängige
		Teilmengen und parallele EDF$^+$-Ausführung ergibt lineare Skalierung.
		Der kritische Pfad (abhängige Tasks) bleibt $O(t_{\max})$.
	\end{proof}
	
	\section{Satz 2: Subgraph-basierte Treiber-Kompatibilität}
	
	\begin{satz}[Treiber-Kompatibilität via Subgraph]\label{satz:drv_compat}
		Sei $G_S = (V_S, E_S)$ der Systemgraph der geladenen Treiber und
		$G_N = (V_N, E_N)$ der Abhängigkeitsgraph eines neuen Treibers.
		Der neue Treiber ist kompatibel genau dann, wenn $G_N \subseteq G_S$.
		Diese Prüfung ist in $O(|V_S|^3)$ entscheidbar.
	\end{satz}
	
	\begin{proof}
		\textbf{Korrektheit der Modellierung:} Kompatibilität erfordert
		$V_N \subseteq V_S$ und eine injektive Abbildung
		$\phi: V_N \to V_S$ mit $\phi(E_N) \subseteq E_S$ --
		genau Subgraph-Isomorphismus.
		
		\textbf{Komplexität:}
		\begin{enumerate}
			\item Signaturberechnung (Definition~\ref{def:sig}) für alle $v \in V_S$: $O(|V_S|)$.
			\item Kandidatenfilterung: nur $v$ mit $\sigma(v) = \sigma(u)$
			werden für $u \in V_N$ geprüft.
			\item Adjazenzprüfung für Kandidatenpaare: $O(|V_N|^2 \cdot |V_S|) = O(n^3)$.
		\end{enumerate}
		
		\textbf{Korrektheit:} Die Signaturfunktion ist notwendige Bedingung
		für Isomorphismus; die Adjazenzprüfung ist hinreichend.
	\end{proof}
	
	\begin{bemerkung}
		Für $n \leq 200$ Treiber läuft die Prüfung in Millisekunden.
		Der $O(n^2)$-Average-Case ermöglicht Echtzeit-Prüfung bei
		dynamisch geladenen Kernel-Modulen.
	\end{bemerkung}
	
	\section{Satz 3: Stochastische Bootzeit-Schranke}
	
	\begin{satz}[Stochastische Bootzeit-Schranke]\label{satz:boottime}
		Seien $C_1, \ldots, C_n$ unabhängige Zufallsvariablen mit
		$C_i \in [a_i, b_i]$. Für die Gesamtbootzeit
		$T_{\mathrm{boot}} = \sum_{i=1}^n C_i$ gilt:
		\[
		\Prob\!\left(T_{\mathrm{boot}} - \E[T_{\mathrm{boot}}] \geq \varepsilon\right)
		\leq \exp\!\left(-\frac{2\varepsilon^2}{\sum_{i=1}^n (b_i - a_i)^2}\right).
		\]
	\end{satz}
	
	\begin{proof}
		Direkte Anwendung der Hoeffding-Ungleichung \cite{Hoeffding1963}:
		Für unabhängige beschränkte Zufallsvariablen $X_i \in [a_i, b_i]$ gilt
		\[
		\Prob\!\left(\sum_i X_i - \sum_i \E[X_i] \geq t\right)
		\leq \exp\!\left(-\frac{2t^2}{\sum_i (b_i-a_i)^2}\right).
		\]
		Mit $X_i = C_i$ und $t = \varepsilon$ folgt die Aussage.
	\end{proof}
	
	\begin{beispiel}
		Für $n = 20$ Komponenten mit $b_i - a_i = 10\,\text{ms}$ und
		$\E[T_{\mathrm{boot}}] = 3\,\text{s}$:
		\[
		\Prob(T_{\mathrm{boot}} > 3{,}1\,\text{s})
		\leq \exp\!\left(-\frac{2 \cdot (0{,}1)^2}{20 \cdot (0{,}01)^2}\right)
		= e^{-10} \approx 4{,}5 \times 10^{-5}.
		\]
		Mit mehr als $99{,}99\,\%$ Wahrscheinlichkeit dauert der Boot unter 3,1\,s.
	\end{beispiel}
	
	\section{Satz 4: Lyapunov-Stabilität des Betriebssystems}
	
	\begin{satz}[OS-Lyapunov-Stabilität]\label{satz:os_lyapunov}
		Das BS/EDF$^+$-System mit Zustandsmatrix $A$ ist asymptotisch
		stabil unter EDF$^+$-Scheduling genau dann, wenn
		$\rho(A) = \max_i |\lambda_i(A)| < 1$,
		bzw.\ $\mathrm{Re}(\lambda_i) < 0$ im kontinuierlichen Analogon.
	\end{satz}
	
	\begin{proof}
		\textbf{Existenz der Lyapunov-Funktion:} Sei $Q = I_4$. Die diskrete
		Lyapunov-Gleichung $A^T P A - P = -Q$ hat eine eindeutige positiv
		definite Lösung $P \succ 0$ genau dann, wenn $\rho(A) < 1$
		\cite{Khalil2002}.
		
		\textbf{V ist Lyapunov-Funktion:} Mit $V(\mathbf{x}) = \mathbf{x}^T P \mathbf{x}$:
		\[
		V(A\mathbf{x}) - V(\mathbf{x})
		= \mathbf{x}^T(A^T P A - P)\mathbf{x}
		= -\mathbf{x}^T Q \mathbf{x}
		= -\|\mathbf{x}\|^2 < 0
		\]
		für $\mathbf{x} \neq \mathbf{0}$.
		
		\textbf{EDF$^+$ als stabilisierendes Scheduling:} Das Scheduling-Gesetz
		$\mathbf{u}_t = K \mathbf{x}_t$ (EDF$^+$, wobei $K$ aus der
		Deadline-Priorität abgeleitet ist) garantiert, dass $A + BK$
		Stabilität erbt, da $(A, B)$ steuerbar durch ACPI/DVFS-Aktuatoren.
	\end{proof}
	
	\begin{korollar}[KS-IV-Vermeidung]
		Falls $\rho(A) < 1$, existiert $\varepsilon > 0$ mit
		$\|\mathbf{x}_t\|_\infty < 1 - \varepsilon$ für alle $t \geq t_0$:
		der KS-IV-Ãœberlastzustand wird asymptotisch vermieden.
	\end{korollar}
	
	\section{Korollar 5: Optimale DVFS-Kurve beim Booten}
	
	\begin{korollar}[IoFET-DVFS Energieoptimalität]\label{kor:dvfs_opt}
		Die IoFET-Sigmoid-Kurve $V^*(f)$ ist die eindeutig energieoptimale
		Spannungs-Frequenz-Kurve für die Boot-Sequenz unter Einhaltung aller
		EDF$^+$-Deadlines.
	\end{korollar}
	
	\begin{proof}
		\textbf{Optimierungsproblem:} Minimiere
		\[
		E_{\mathrm{boot}} = \sum_{i=1}^{n} \alpha_i C_L V_i^2 f_i C_i
		\]
		unter $f_i C_i \leq D_i$ und $V_{\min} \leq V_i \leq V_{\max}$.
		
		\textbf{KKT-Bedingungen:} Die stationäre Bedingung
		$\partial \mathcal{L}/\partial V_i = 2\alpha_i C_L V_i f_i C_i + \lambda_{V,i} = 0$
		liefert $V_i^* \propto 1/\sqrt{f_i}$, was einer sigmoiden Kurve
		entspricht, wenn die Frequenzabhängigkeit des IoFET-Modells berücksichtigt wird.
		
		\textbf{Eindeutigkeit:} Da $E(f)$ nach Satz~\ref{satz:dvfs_convex}
		konvex ist, ist das Minimum eindeutig. Kein anderes Steuergesetz
		erzielt geringere Boot-Energie.
	\end{proof}
	
	\section{Satz 6: PQC-sichere Speicheradressierung}
	
	\begin{satz}[Quantensichere Speicheradressierung]\label{satz:pqc_mem}
		Sei $\mathbf{s} \in \mathbb{Z}_q^n$ ein geheimer Schlüssel und
		$\Phi: \mathbb{Z}^{64} \to \mathbb{Z}_q^n$ eine injektive Hash-Funktion.
		Das LWE-Adressierungsschema
		\[
		\mathrm{PA}(a) = \bigl(\mathbf{a},\; \langle \mathbf{a}, \mathbf{s} \rangle
		+ e + \Phi(a) \pmod{q}\bigr)
		\]
		ist quantensicher: Kein PPT-Quantenangreifer kann $\mathrm{PA}(a) \mapsto a$
		mit nicht-vernachlässigbarer Wahrscheinlichkeit bestimmen, sofern LWE
		quantenschwer ist.
	\end{satz}
	
	\begin{proof}
		\textbf{Sicherheitsreduktion:} Angenommen, $\mathcal{A}_Q$ bestimmt
		$\mathrm{PA}(a) \mapsto a$ mit Vorteil $\varepsilon > 0$. Dann kann
		$\mathcal{A}_Q$ genutzt werden, um das LWE-Problem zu lösen: Gegeben
		LWE-Sample $(\mathbf{a}, b)$, setze $b' = b - \Phi(0)$.
		$\mathcal{A}_Q$ müsste $\mathbf{s}$ aus $(\mathbf{a}, b')$ ableiten --
		das ist das LWE-Problem.
		
		\textbf{Quantenschwere von LWE:} Regev \cite{Regev2009} zeigt, dass
		LWE mindestens so schwer ist wie $\mathrm{SIVP}_\gamma$ im Worst Case.
		Es ist kein effizienter Quantenalgorithmus für $\mathrm{SIVP}$ bekannt.
		
		\textbf{Effizienz:} LWE-Entschlüsselung ($n = 1024$, $q = 12289$)
		benötigt $< 1\,\mu\text{s}$ auf moderner Hardware. BMM \cite{Epp2026boolmm}
		optimiert den TLB-Cache auf $O(n^{2.37})$.
	\end{proof}
	
	\begin{bemerkung}
		Die vollständige LWE-Adressierung aller Pages ist für
		Hochsicherheitsanwendungen vorgesehen. Für Standard-Betrieb genügt
		selektive Anwendung auf Kernel-Speicherbereiche (Page Tables, IDT, APIC).
	\end{bemerkung}
	
	% ==============================================================
	%  KAPITEL 15 -- ZUSAMMENFASSUNG UND AUSBLICK
	%  Erweiterter Ausblick -- Ersatz für den bisherigen Abschnitt
	%  "Offene Fragen und Ausblick" in bs.tex
	% ==============================================================
	
	\chapter{Zusammenfassung und Ausblick}
	
	\section{Zusammenfassung}
	
	Diese Arbeit hat BS/EDF$^+$ als formal fundiertes Betriebssystem-Konzept
	vorgestellt, das EDF$^+$-Scheduling als einheitliches Paradigma von der
	Hardware-Initialisierung bis zur Benutzeranwendung einsetzt.
	
	Die sechs neuen Hauptergebnisse:
	\begin{enumerate}
		\item \textbf{Satz~14.1:} Boot-Sequenz als EDF$^+$-Task-System --
		formale Korrektheit der Initialisierungsreihenfolge.
		\item \textbf{Satz~14.2:} Treiber-Kompatibilität via Subgraph-Isomorphismus
		in $O(n^3)$.
		\item \textbf{Satz~14.3:} Hoeffding-Schranke für die Bootzeit:
		$> 99{,}99\,\%$ Wahrscheinlichkeit unter $t_{\max}$.
		\item \textbf{Satz~14.4:} Lyapunov-Stabilität unter EDF$^+$ vermeidet
		KS-IV-Überlastzustände.
		\item \textbf{Korollar~14.5:} IoFET-Sigmoid ist energieoptimale
		DVFS-Kurve beim Booten.
		\item \textbf{Satz~14.6:} LWE-basierte Speicheradressierung ist
		Post-Quanten-sicher.
	\end{enumerate}
	
	% ==============================================================
	\section{Offene Fragen und erweiterter Ausblick}
	% ==============================================================
	
	Der in dieser Arbeit entwickelte BS/EDF$^+$-Rahmen bildet eine wissenschaftliche
	Grundlage, die in mehrere Richtungen erheblich vertieft und ausgebaut werden kann.
	Der folgende Ausblick gliedert sich in kurzfristige (1--3 Jahre),
	mittelfristige (3--7 Jahre) und langfristige (7--15 Jahre) Forschungsziele
	sowie in übergreifende Themen zur gesellschaftlichen und industriellen
	Wirkung des hier entwickelten Ansatzes.
	
	% ---------------------------------------------------------------
	\subsection{Formale Verifikation: Von Papierbeweisen zu maschinengeprüften Zertifikaten}
	% ---------------------------------------------------------------
	
	Die sechs Sätze in Kapitel~14 sind vollständige mathematische Beweise auf
	Papierebene. Der unmittelbarste und wichtigste nächste Schritt ist ihre
	\emph{maschinelle Verifikation} in einem interaktiven Beweisassistenten.
	
	\subsubsection{Isabelle/HOL und Coq}
	
	Isabelle/HOL \cite{Nipkow2002} eignet sich besonders für Satz~14.1 und
	Satz~14.2, da beide auf endlichen Graphenstrukturen operieren, für die
	reiche HOL-Bibliotheken (z.\,B.\ \texttt{Graph\_Theory}) existieren.
	Der Beweis von Satz~14.1 lässt sich in Isabelle als parametrisierter
	Satz über beliebige EDF$^+$-Task-Systeme mit $n$ Phasen formulieren:
	
	\[
	\forall\, n \in \mathbb{N},\; \forall\,\tau_1,\ldots,\tau_n.\;
	U_{\mathrm{EDF}^+} \le 1 \;\Longrightarrow\;
	\text{sched}(\tau_1,\ldots,\tau_n) \text{ ist terminierend}
	\]
	
	Coq bietet darüber hinaus die Möglichkeit, aus dem verifizierten Beweis
	direkt \emph{zertifizierten C-Code} zu extrahieren (via \texttt{Extraction}),
	der in den BS/EDF$^+$-Kernel integriert werden kann -- ein Ansatz, der
	bereits im seL4-Projekt \cite{Klein2009} für den Kernel-Korrektheitsbeweis
	erfolgreich eingesetzt wurde.
	
	\subsubsection{TLA\texorpdfstring{$^+$}{+} für concurrent Boot-Sequenzen}
	
	TLA$^+$ \cite{Lamport2002} ist prädestiniert für die Modellierung
	nebenläufiger Boot-Phasen auf PYMCA-8-Mehrkern\-systemen. Insbesondere
	die Synchronisationsinvarianten zwischen den acht Kernen (Kapitel~6) lassen
	sich als TLA$^+$-Spezifikation ausdrücken und mit dem TLC-Model-Checker
	auf Deadlock-Freiheit und Liveness-Eigenschaften prüfen. Eine vollständige
	TLA$^+$-Spezifikation von BS/EDF$^+$ könnte als maschinengeprüftes
	Referenzdokument für zukünftige Kernel-Implementierungen dienen.
	
	\subsubsection{Parametrische Timed Automata}
	
	Der unendliche Zustandsraum der Boot-Sequenz (Satz~14.3) ist ein klassisches
	Problem für \emph{Parametrische Timed Automata} (PTA, Alur et al.\ 1993).
	Mit Werkzeugen wie IMITATOR lassen sich parametrische Zeitschranken
	$t_{\max}(\Delta, n)$ als Funktion der Systemparameter automatisch synthesieren --
	eine direkte Erweiterung der in Satz~14.3 manuell bewiesenen
	Hoeffding-Schranke zu einem algorithmisch berechenbaren Zertifikat.
	
	% ---------------------------------------------------------------
	\subsection{Portierung auf ARM, RISC-V und heterogene Architekturen}
	% ---------------------------------------------------------------
	
	BS/EDF$^+$ wurde für x86-64 entwickelt, ist aber konzeptionell
	ISA-unabhängig. Drei Portierungsziele erscheinen besonders vielversprechend.
	
	\subsubsection{ARM Cortex-A78 und die mobile Domäne}
	
	Der ARM Cortex-A78 ist der Kern moderner Smartphone-SoCs
	(Qualcomm Snapdragon 888, Samsung Exynos 2200). Die DVFS-Implementierung
	muss von x86-64 MSRs auf die ARM \emph{SCMI} (System Control and
	Management Interface) portiert werden; die IoFET-Sigmoid-Kurve (Korollar~14.5)
	ist architekturunabhängig und kann direkt übernommen werden.
	Für den mobilen Einsatz ist insbesondere die Integration mit
	Android-Treiber-Stacks (Binder IPC, HIDL) durch Subgraph-Isomorphismus
	(Satz~14.2) interessant: Das Framework \cite{Epp2026mobde} liefert
	hierfür bereits Vorarbeiten.
	
	\subsubsection{RISC-V RV64GC: offene ISA-Verifikation}
	
	RISC-V (RV64GC) ermöglicht es, die gesamte ISA-Spezifikation quelloffen
	zu halten. In Kombination mit formaler Verifikation in Isabelle/HOL
	entsteht damit ein durchgehend verifizierbarer Stack: von der
	Transistorschicht (IoFET, \cite{Epp2026iono}) über die ISA bis zum
	OS-Scheduler. Das Ziel ist ein \emph{durchgehend beweisbares System}
	ohne vertrauenswürdige Komponenten -- ein \emph{Trusted Computing Base}
	von theoretisch null Elementen, sofern die Hardware-Verifikation
	abgeschlossen ist.
	
	\subsubsection{Heterogene CPU-GPU-NPU-Architekturen}
	
	Moderne Systeme on Chip (NVIDIA Orin, Apple M3) kombinieren CPU-Kerne,
	Grafik\-prozessoren und Neural Processing Units auf einem Die.
	BS/EDF$^+$ müsste auf einen \emph{Heterogeneous EDF$^+$}-Algorithmus
	(HEDF$^+$) erweitert werden, der Tasks dynamisch auf die am besten
	geeignete Ausführungseinheit zuweist. Die Subgraph-Kompatibilitätsprüfung
	(Satz~14.2) bietet einen natürlichen Rahmen: Ein Task-Graph wird mit
	dem Hardware-Ressourcengraph verglichen; der Isomorphismus liefert
	die optimale Zuordnung. Die Laufzeit bleibt $O(n^3)$ auf der
	kombinierten Graphstruktur.
	
	% ---------------------------------------------------------------
	\subsection{KI-gestützte Systemoptimierung}
	% ---------------------------------------------------------------
	
	Maschinelles Lernen bietet erhebliche Möglichkeiten, die in dieser Arbeit
	entwickelten formalen Schranken empirisch zu verschärfen.
	
	\subsubsection{LSTM-basierte Bootzeit-Vorhersage}
	
	Die stochastische Schranke in Satz~14.3 setzt voraus, dass die
	Ausführungszeiten $C_i$ der Boot-Tasks bekannt sind. In der Praxis
	variieren diese durch Hardware-Alterung, Temperatureffekte und
	Speicherfehler. Ein LSTM-Netzwerk, das auf Millionen von Boot-Traces
	trainiert wird, kann $C_i$ für den nächsten Boot mit weniger als
	$1\,\%$ relativer Abweichung vorhersagen und damit die Hoeffding-Schranke
	adaptiv anpassen. Das folgende Trainings-Ziel formalisiert dies:
	
	\[
	\min_{\Theta}\; \mathbb{E}\!\left[\sum_{i=1}^{n}
	\left(\hat{C}_i^{\Theta} - C_i\right)^2\right],
	\]
	
	wobei $\Theta$ die LSTM-Parameter und $\hat{C}_i^{\Theta}$ die vorhergesagte
	Ausführungszeit sind. Mit einem kalibrierten $\hat{C}_i$ kann
	$t_{\max}$ um typischerweise $15$--$25\,\%$ reduziert werden, ohne die
	Verletzungswahrscheinlichkeit zu erhöhen.
	
	\subsubsection{Reinforcement Learning für adaptives DVFS}
	
	Die IoFET-Sigmoid-DVFS-Kurve (Korollar~14.5) ist statisch optimal
	unter der Annahme eines bekannten Arbeitslastprofils. Ein Reinforcement-Learning-Agent
	(z.\,B.\ PPO oder SAC) könnte die DVFS-Parameter in Echtzeit anpassen,
	indem er die thermische Zustandsmatrix $A$ des Systems (Satz~14.4)
	als Beobachtungsraum und den Energieverbrauch als negativen Reward verwendet.
	Das Lyapunov-Stabilitätskriterium aus Satz~14.4 kann dabei als
	harte Constraint in den RL-Optimierungsprozess eingebettet werden
	(\emph{Constrained RL}), sodass der Agent niemals den KS-IV-Ãœberlastzustand
	induziert.
	
	\subsubsection{Anomalieerkennung im Kernel}
	
	Der Subgraph Algorithmus \cite{Epp2026subgraph} ermöglicht eine neuartige
	Form der Kernel-Integritätsprü\-fung: Der Graph der aktiven Kernel-Subsysteme
	(Treiber, Interrupt-Handler, IPC-Kanäle) wird kontinuierlich als
	\emph{Graph-Profil} gespeichert und mit einem Referenzprofil verglichen.
	Eine Abweichung, die sich nicht als Subgraph-Isomorphismus darstellen lässt,
	signalisiert eine anomale Systemveränderung -- z.\,B.\ durch Rootkits oder
	Speicherkorrumption. In Kombination mit LWE-gesicherter Speicheradressierung
	(Satz~14.6) und der stochastischen Archiv-Lösung \cite{Epp2026archv}
	entsteht ein dreistufiges Integritätssystem:
	(1) kryptographische Speichersicherung,
	(2) graphenbasierte Strukturintegrität,
	(3) stochastische Archivintegrität.
	
	% ---------------------------------------------------------------
	\subsection{Quantencomputer-Resilienz und Post-Quanten-Migration}
	% ---------------------------------------------------------------
	
	\subsubsection{Empirische Härtungsversuche mit Gitterreduktionsalgorithmen}
	
	Sobald Quantencomputer mit mehr als 1000 \emph{logischen} Qubits verfügbar
	sind (nach aktuellem Stand der Technik voraussichtlich ab 2030--2035),
	sollte die Sicherheit des LWE-Schemas (Satz~14.6) mit dem BKZ-2.0-Algorithmus
	\cite{Chen2011} unter realen Quantenfehlerkorrekturoverheads überprüft werden.
	Für $n = 1024$, $q = 12289$ liegt die klassische Angriffskomplexität bei
	$2^{218}$ Operationen; mit Quantenbeschleunigung (Grover-Vorfaktor
	$\sqrt{\cdot}$) sinkt sie auf $\approx 2^{109}$ -- immer noch weit jenseits
	jeder praktischen Angreifbarkeit bei korrekter Parameterwahl.
	
	\subsubsection{Migration zu CRYSTALS-Kyber und CRYSTALS-Dilithium}
	
	Der NIST-Post-Quanten-Standard hat im Jahr 2024 CRYSTALS-Kyber \cite{Bos2018}
	für Schlüsselaustausch und CRYSTALS-Dilithium \cite{Ducas2018} für
	digitale Signaturen standardisiert. BS/EDF$^+$ sollte beide Schemata
	als native Kernel-Primitive integrieren:
	
	\begin{itemize}
		\item \textbf{Kyber-1024} als Ersatz für RSA-2048 im Secure-Boot-Prozess:
		Der Bootloader verifiziert die Kernel-Signatur mit Kyber-1024,
		das gegen Shors Algorithmus \cite{Shor1994} immun ist.
		\item \textbf{Dilithium-3} für die Kernel-Modul-Signierung:
		Jedes Treibermodul trägt eine Dili\-thium-Signatur; der Subgraph
		Algorithmus prüft sowohl die kryptographische Integrität als auch
		die strukturelle Kompatibilität (Satz~14.2).
	\end{itemize}
	
	\subsubsection{Hybride klassisch-quantensichere Adressierung}
	
	Satz~14.6 beschreibt eine vollständige LWE-Adressierung, die für
	Hochsicherheitsanwendungen geeignet ist. Für den Standardbetrieb
	ist eine \emph{hybride Strategie} effizienter: Konventionelle
	virtuelle Adressierung (x86-64 4-Ebenen-Paging) schützt reguläre
	Benutzer-Pages; LWE-Adressierung wird selektiv auf
	sicherheitskritische Kernel-Regionen (Page Tables, IDT, APIC-Register)
	angewendet. Diese Hybrid-Strategie reduziert den Overhead von
	$O(n^{2.37})$ auf $O(k \cdot n^{2.37}/n)$, wobei $k$ die Anzahl
	der geschützten Pages ist ($k \ll n$). Eine formale
	Analyse dieses Trade-offs ist Gegenstand zukünftiger Arbeit.
	
	% ---------------------------------------------------------------
	\subsection{Erweiterung des PYMCA-8-Modells auf viele Kerne und NUMA}
	% ---------------------------------------------------------------
	
	Die PYMCA-8-Architektur (Kapitel~6, \cite{Epp2026pymca8}) beschreibt
	acht Kerne mit Python-Speicher\-modell und IoFET-DVFS. Die natürliche
	Erweiterung führt in zwei Richtungen.
	
	\subsubsection{PYMCA-$n$: Skalierung auf beliebig viele Kerne}
	
	Eine theoretische Verallgemeinerung auf $n$ Kerne erfordert:
	
	\begin{enumerate}
		\item Eine formale Beweis-Erweiterung von Satz~14.4 (Lyapunov-Stabilität)
		auf $n \times n$-Systemmatrizen $A \in \mathbb{R}^{n \times n}$, wobei
		NUMA-Latenzeffekte als Off-Diagonal-Einträge modelliert werden.
		\item Eine skalierbare Variante des Subgraph Algorithmus
		\cite{Epp2026subgraph} für $n$-Kern-Treibergraphen, die die
		$O(n^3)$-Schranke durch parallele Subgraph-Zerlegung auf
		$O(n^2 \log n)$ reduziert.
		\item Eine Erweiterung des stochastischen Archiv-Problems
		\cite{Epp2026archv} auf verteilte Speicher mit NUMA-Latenzmodell.
	\end{enumerate}
	
	\subsubsection{Integration mit AUTOSAR-Adaptive auf Mehrkernsystemen}
	
	AUTOSAR Adaptive Platform (AP, R23-11 \cite{AUTOSAR2023}) zielt auf
	hochleistungsfähige SoCs mit dynamischer Dienstarchitektur.
	BS/EDF$^+$ könnte als Execution Platform für AUTOSAR-AP dienen:
	Der EDF$^+$-Scheduler übernimmt das Task-Scheduling der
	AUTOSAR-Execution-Management-Schicht; der Subgraph Algorithmus
	prüft Service-Kompatibilität im AUTO\-SAR-Dienstverzeichnis.
	Damit entsteht eine formale Brücke zwischen dem akademischen
	BS/EDF$^+$-Konzept und dem industriellen AUTOSAR-Standard --
	ein wesentlicher Schritt zur praktischen Verwertung dieser Arbeit.
	
	\subsubsection{EtherCAT-EDF\texorpdfstring{$^+$}{+}-Standardisierung}
	
	Die Integration von EtherCAT \cite{IEC61158} in das EHAE-Framework
	\cite{Epp2026ethercate} mit EDF$^+$-Scheduling sollte bei der
	EtherCAT Technology Group (ETG) als Standarderweiterung vorgeschlagen werden.
	Konkret ermöglicht EDF$^+$ eine \emph{dynamische Slot-Vergabe} im
	EtherCAT-Frame: Statt statisch konfigurierten Zeitfenstern (wie in
	der aktuellen ETG-Spezifikation) vergibt EDF$^+$ Ãœbertragungsslots
	nach Deadline-Priorität. Satz~14.1 garantiert die Schedulierbarkeit
	unter der Bedingung $U_{\mathrm{EDF}^+} \le 1$; für typische
	Industrieinstallationen mit $n = 128$ EtherCAT-Slaves liegt
	$U$ bei $\approx 0{,}72$ -- hinreichend weit unterhalb der
	Stabilitätsgrenze.
	
	% ---------------------------------------------------------------
	\subsection{Energiemanagement: IoFET-DVFS und Li-Ionen-Integration}
	% ---------------------------------------------------------------
	
	\subsubsection{Dynamische Temperaturmodelle für IoFET-Transistoren}
	
	Korollar~14.5 beweist die Optimalität der IoFET-Sigmoid-DVFS-Kurve
	unter der Annahme konstanter Temperatur. Reale Systeme unterliegen
	jedoch thermischen Transienten: Die Drain-Strom-Charakteristik
	$I_D = I_0 \sinh\!\left(\frac{V_{GS}}{nV_T}\right)$ ist stark temperaturabhängig
	($V_T = k_BT/q$). Eine Erweiterung des Modells \cite{Epp2026iono}
	auf temperaturabhängige IoFET-Parameter führt zu einem
	\emph{thermoelektrisch gekoppelten DVFS-Optimierungsproblem}:
	
	\[
	\min_{f(t),\, V(t)}\; \int_0^T P(f,V,T(t))\,\mathrm{d}t
	\quad\text{s.\,t.}\quad
	\dot{T} = \alpha\, P(f,V,T) - \beta\,(T - T_{\mathrm{amb}}),
	\]
	
	wobei $T(t)$ die Chip-Temperatur ist und $\alpha, \beta$ thermische
	Transferkoeffizienten bezeichnen. Dieses Problem ist ein Optimal-Control-Problem
	und kann mit Pontrjaginschen Maximumprinzip oder numerischer
	Kollokation gelöst werden.
	
	\subsubsection{Hybride IoFET-Li-Ionen-Energiepolitik}
	
	Die KI-gestützte Akkuverwaltung \cite{Epp2026liionp} und das
	IoFET-DVFS-Modell \cite{Epp2026iono} wurden bisher separat entwickelt.
	Eine integrierte Energiepolitik, die den Akkuzustand (State-of-Charge,
	State-of-Health) als Eingabe des DVFS-Reglers verwendet, könnte
	die Gesamtsystemlebensdauer signifikant verlängern. Formal lässt sich
	dies als \emph{Markov-Decision-Process} (MDP) formulieren:
	
	\[
	\max_\pi\; \mathbb{E}_\pi\!\left[\sum_{t=0}^{\infty}
	\gamma^t R(s_t, a_t)\right],
	\]
	
	wobei der Zustandsraum $s_t = (f_t, V_t, \mathrm{SoC}_t, T_t)$
	und die Aktion $a_t$ die Wahl der DVFS-Stufe ist.
	Der Reward $R$ balanciert Performanz und Akkuverschleiß.
	
	% ---------------------------------------------------------------
	\subsection{Dateisystem und Speicherverwaltung}
	% ---------------------------------------------------------------
	
	\subsubsection{Subgraph-basierte Deduplizierung auf NVMe-Speicher}
	
	Das Dateisystem-Deduplizierungsframework \cite{Epp2026dusgraph} nutzt
	Subgraph-Isomorphismus zur Erkennung strukturell ähnlicher Datenblöcke.
	Aktuelle Implementierungen operieren auf Block-Ebene; eine Erweiterung
	auf \emph{semantische Deduplizierung} -- d.\,h.\ inhaltsbezogene Ähnlichkeit
	auf Ebene von Dokumentstrukturen, Code-Abstrakt-Syntaxbäumen (AST) oder
	Vektordatenbanken (VADIS \cite{Epp2026vadis}) -- würde erhebliche
	Speichereinsparungen ermöglichen. Für NVMe-SSDs mit $10^9$ Datenobjekten
	ist eine parallele Subgraph-Verarbeitung auf PYMCA-8 mit $O(n^2\log n)$
	Komplexität erforderlich.
	
	\subsubsection{Persistente Speicherklassen (NVDIMM, CXL)}
	
	Compute Express Link (CXL) 3.0 ermöglicht Memory-Pooling zwischen
	mehreren Hosts. BS/EDF$^+$ sollte CXL als neue Speicherklasse
	modellieren: CXL-Speicher hat höhere Latenz als DRAM, aber niedrigere
	als NVMe; das stochastische Archiv-Modell \cite{Epp2026archv} muss
	um eine dritte Speicherebene (DRAM / CXL / NVMe) erweitert werden.
	Die optimale Batch-Timer-Strategie ändert sich dabei qualitativ:
	Für CXL-Zugriffe mit Poisson-Ankünften der Rate $\lambda_{CXL}$
	ergibt sich ein modifiziertes Optimierungsproblem für den
	Erwartungswert der Latenz.
	
	% ---------------------------------------------------------------
	\subsection{Graphentheoretische Systemverifikation}
	% ---------------------------------------------------------------
	
	\subsubsection{HS-DFS-Erweiterung für dynamische Kernel-Graphen}
	
	Der Subgraph Algorithmus \cite{Epp2026subgraph} und der
	HS-DFS-Algorithmus (Graph-of-Graphs, \cite{Epp2026msubgraph}) operieren
	aktuell auf statischen Graphen. Echtzeit-Kernel-Systeme sind jedoch
	dynamisch: Treiber werden geladen und entladen, IPC-Verbindungen
	werden aufgebaut und beendet. Eine Erweiterung auf \emph{dynamische
		Graph-Ströme} (\emph{graph streams}) mit Online-Subgraph-Erkennung
	in amortisierter $O(\log n)$-Zeit pro Änderungsoperation ist ein
	wichtiger offener Forschungspunkt.
	
	\subsubsection{Boolesche Matrixmultiplikation für Cache-Kohärenz}
	
	Die Boolesche Matrixmultiplikation in $O(n^{2.37})$
	\cite{Epp2026boolmm} wurde bisher für TLB-Cache-Opti\-mierung
	genutzt. Eine direkte Anwendung liegt in der \emph{Cache-Kohärenz-Analyse}
	auf PYMCA-8: Die Abhängigkeitsmatrix $D \in \{0,1\}^{n \times n}$
	zwischen $n$ Tasks wird als Boolesches Matrix-Produkt berechnet;
	Einträge $D_{ij} = 1$ signalisieren potenzielle Cache-Invalidierungen.
	Die schnelle BMM ermöglicht es, diese Analyse in jedem
	Scheduling-Zyklus ($\approx 1\,\mu\text{s}$) neu auszuführen.
	
	% ---------------------------------------------------------------
	\subsection{Sicherheitsarchitektur und Isolation}
	% ---------------------------------------------------------------
	
	\subsubsection{Capability-basierte Zugriffskontrolle}
	
	Das aktuelle BS/EDF$^+$-Sicherheitsmodell basiert auf LWE-Adressierung
	(Satz~14.6) und Subgraph-Integritätsprüfung. Eine Erweiterung um
	\emph{Capabilities} (im Sinne von seL4 oder CHERI) würde eine
	fein\-granulare Zugriffskontrolle auf Objekt\-ebene ermöglichen:
	Jede Capability $c = (o, p, \sigma)$ kodiert ein Objekt $o$,
	eine Permissionsmaske $p$ und eine kryptographische Signatur $\sigma$
	(Dilithium-3). Der Subgraph Algorithmus prüft, ob der Capability-Graph
	des anfragenden Tasks ein gültiger Subgraph des Ressourcen-Graphen ist --
	eine natürliche Erweiterung von Satz~14.2 auf Sicherheitspolicies.
	
	\subsubsection{Isolierte Ausführungsumgebungen für AUTOSAR-Partitionen}
	
	ISO~26262 \cite{ISO26262} fordert räumliche und zeitliche
	Isolation zwischen AUTOSAR-Partitionen (ASIL-A bis ASIL-D).
	BS/EDF$^+$ kann diese Isolation durch eine Kombination aus
	(1) EDF$^+$-Time-Partitioning (zeitliche Isolation),
	(2) LWE-gesicherter Speicheradressierung per Partition (räumliche Isolation)
	und (3) Subgraph-Kompatibilitätsprüfung vor jedem Treiberladen gewährleisten.
	Eine formale Verifikation dieser Drei-Ebenen-Isolation in Isabelle/HOL
	würde eine ASIL-D-Zertifizierung von BS/EDF$^+$ ermöglichen --
	eine erhebliche industrielle Wirkung.
	
	% ---------------------------------------------------------------
	\subsection{Systemzustandsklassifikation und Echtzeit-Diagnostik}
	% ---------------------------------------------------------------
	
	\subsubsection{Erweiterung der KS-Klassifikation auf fünf Ebenen}
	
	Das aktuelle Modell \cite{Epp2026sysstate} unterscheidet vier
	Systemzustände KS-I (stabil) bis KS-IV (überlastet). Eine Erweiterung
	auf fünf Ebenen durch Hinzunahme eines \emph{KS-V: Degradierter
		Betrieb} würde qualitativ neue Reaktionsstrategien ermöglichen:
	Im KS-V-Zustand reduziert der Kernel proaktiv die Anzahl aktiver
	Tasks auf eine sichere Teilmenge, die durch Lyapunov-Stabilität
	(Satz~14.4) garantiert stabil bleibt. Der Ãœbergang KS-IV $\to$
	KS-V wäre durch eine Hysterese-Bedingung geregelt, um Flatter-Effekte
	zu vermeiden.
	
	\subsubsection{Echtzeit-Spektralanalyse der Systemzustands\-matrix}
	
	Satz~14.4 zeigt, dass Lyapunov-Stabilität äquivalent zu
	$\max_i \mathrm{Re}(\lambda_i(A)) < 0$ ist. Eine kontinuierliche
	Online-Ãœberwachung des Spektralradius $\rho(A)$ in Echtzeit erfordert
	effiziente inkrementelle Eigenwertberechnung. Mit der
	Borst-Golub-Reinsch-Methode für \emph{rank-1-updates} der
	Systemmatrix $A \mapsto A + uv^T$ lassen sich
	Eigenwertänderungen in $O(n^2)$ statt $O(n^3)$ berechnen --
	ein wesentlicher Effizienzgewinn für den $n = 8$ Kern-Fall
	von PYMCA-8.
	
	% ---------------------------------------------------------------
	\subsection{Langfristige Vision: BS/EDF\texorpdfstring{$^+$}{+} als offenes Ökosystem}
	% ---------------------------------------------------------------
	
	\subsubsection{Open-Source-Kernel-Implementierung}
	
	Alle in dieser Arbeit beschriebenen Algorithmen und Konzepte sind
	implementierungsoffen und unter MIT-Lizenz auf
	\texttt{github.com/hjstephan86/science/tree/main/science/} veröffentlicht. Die nächste Stufe
	ist ein ausführbarer Proof-of-Concept-Kernel in Rust:
	
	\begin{enumerate}
		\item \textbf{Bootloader:} 16-Bit-Realmode $\to$ 32-Bit-Protected-Mode
		$\to$ 64-Bit-Long-Mode, verifiziert in TLA$^+$.
		\item \textbf{EDF$^+$-Scheduler:} Rust-\texttt{no\_std}-Implementierung
		mit formaler Korrektheit durch Coq-Extraktion.
		\item \textbf{Subgraph-Modul:} $O(n^3)$-Treiber-Kompatibilitätsprüfung
		als Kernel-Subsystem.
		\item \textbf{LWE-Speichertreiber:} CRYSTALS-Kyber-basierter
		Page-Table-Walk für Hochsicherheitsseiten.
		\item \textbf{IoFET-DVFS-Regler:} Sigmoid-Kurvenparametrisierung
		mit Online-Temperatur\-korrektur.
	\end{enumerate}
	
	\subsubsection{Interoperabilität mit Linux-Treibern via AUTOSAR-HAL}
	
	Eine vollständige Neuentwicklung aller Gerätetreiber ist
	mittelfristig nicht realistisch. Stattdessen kann BS/EDF$^+$
	über eine \emph{Hardware Abstraction Layer} (HAL) kompatibel
	zu bestehenden Linux-Treibern (via AUTOSAR-HAL-Schicht,
	\cite{Epp2026autsrepy}) gemacht werden: Linux-Treiber werden
	als EDF$^+$-Tasks mit konservativ geschätzten WCET-Werten
	modelliert; Satz~14.2 prüft Kompatibilität automatisch.
	Damit wird der Ãœbergang von Linux zu BS/EDF$^+$ zu einem
	inkrementellen Migrationspfad.
	
	\subsubsection{Referenzimplementierung für die Lehre}
	
	BS/EDF$^+$ ist als Lehrgegenstand in Betriebssystem-Vorlesungen
	besonders geeignet, weil alle Konzepte -- EDF-Scheduling,
	formale Verifikation, Kryptographie, Graphenalgorithmen --
	in einem kohärenten Rahmen zusammengeführt werden.
	Eine \emph{QEMU-basierte Lehrumgebung} mit vorkompilierten
	Kernel-Images und schrittweisem Debugger würde Studierenden
	ermöglichen, jeden der sechs Hauptsätze dieser Arbeit
	direkt am laufenden System zu beobachten.
	
	% ---------------------------------------------------------------
	\subsection{Gesellschaftliche und industrielle Perspektive}
	% ---------------------------------------------------------------
	
	\subsubsection{Sicherheitskritische Infrastruktur}
	
	Mit der zunehmenden Digitalisierung kritischer Infrastrukturen
	(Energienetze, Wasserversorgung, autonome Fahrzeuge) steigt der
	Bedarf an formal verifizierten Betriebssystemen dramatisch.
	BS/EDF$^+$ bietet einen Weg, wie ein vollständig beweisbares
	Echtzeit-OS für ASIL-D-Systeme aussehen könnte -- ein Beitrag,
	der weit über die akademische Forschung hinausgeht.
	
	\subsubsection{Software-Defined Vehicles}
	
	Die AUTOSAR-OTA-Architektur \cite{Epp2026autsre} und 
	EDF$^+$-basierte Treiber-Validierung (Satz~14.2) sind direkte
	Bausteine für Software-Defined Vehicles (SDV): Fahrzeuge, deren
	gesamte Funktionalität durch Software definiert und über-the-air
	aktualisiert wird. BS/EDF$^+$ liefert das formale Fundament für
	die Sicherheitsnachweise, die regulatorische Zulassungsbehörden
	(KBA, UN-ECE WP.29) für SDV-Systeme fordern werden.
	
	\subsubsection{Quantensichere Infrastruktur bis 2035}
	
	Satz~14.6 positioniert BS/EDF$^+$ als frühen Vertreter
	einer Generation von Betriebssystemen, die bereits heute
	gegen zukünftige Quantenangriffe gehärtet sind. Mit der
	bevorstehenden Migration nationaler Sicherheitsbehörden
	auf Post-Quanten-Kryptographie (BSI TR-02102, NIST SP 800-208)
	ist die in dieser Arbeit entwickelte LWE-Speicheradressierung
	ein zeitgemäßer und zukunftssicherer Ansatz.
	
	\begin{bemerkung}
		Die in diesem Ausblick beschriebenen Forschungsziele sind nicht
		als sequenziell zu verstehen. Viele Pfade -- formale Verifikation,
		ARM-Portierung, RL-basiertes DVFS, Post-Quanten-Migration --
		sind parallel verfolgbar und gegenseitig bereichernd.
		BS/EDF$^+$ ist kein abgeschlossenes System, sondern ein
		lebendiges Forschungsprogramm, dessen innere Konsistenz durch
		die sechs Hauptsätze dieser Arbeit garantiert wird.
	\end{bemerkung}
	
	% ==============================================================
	%  LITERATURVERZEICHNIS
	% ==============================================================
	\begin{thebibliography}{99}
		
		\bibitem{Liu1973}
		C.\,L.\ Liu und J.\,W.\ Layland,
		``Scheduling algorithms for multiprogramming in a hard-real-time environment,''
		\textit{Journal of the ACM}, Bd.~20, Nr.~1, S.~46--61, 1973.
		
		\bibitem{Dertouzos1974}
		M.\,L.\ Dertouzos,
		``Control robotics: The procedural control of physical processes,''
		in \textit{Proc.\ IFIP Congress}, S.~807--813, 1974.
		
		\bibitem{Hoeffding1963}
		W.\ Hoeffding,
		``Probability inequalities for sums of bounded random variables,''
		\textit{Journal of the American Statistical Association}, Bd.~58, Nr.~301,
		S.~13--30, 1963.
		
		\bibitem{Cormen2022}
		T.\,H.\ Cormen, C.\,E.\ Leiserson, R.\,L.\ Rivest und C.\ Stein,
		\textit{Introduction to Algorithms},
		4.~Aufl.\ Cambridge, MA: MIT Press, 2022.
		
		\bibitem{Buttazzo2011}
		G.\,C.\ Buttazzo,
		\textit{Hard Real-Time Computing Systems},
		3.~Aufl.\ New York: Springer, 2011.
		
		\bibitem{AUTOSAR2023}
		AUTOSAR Consortium,
		\textit{AUTOSAR Classic Platform -- Specification of Operating System},
		Release~R23-11, 2023.
		
		\bibitem{ISO26262}
		ISO,
		\textit{ISO~26262: Road Vehicles -- Functional Safety},
		2.~Aufl., 2018.
		
		\bibitem{Intel2023}
		Intel Corporation,
		\textit{Intel\textregistered{} 64 and IA-32 Architectures Software
			Developer's Manual}, Vol.~1--3, 2023.
		
		\bibitem{IEC61158}
		IEC,
		\textit{IEC~61158: Industrial Communication Networks -- Fieldbus
			Specifications}, 2019.
		
		\bibitem{Regev2009}
		O.\ Regev,
		``On lattices, learning with errors, random linear codes, and cryptography,''
		\textit{Journal of the ACM}, Bd.~56, Nr.~6, S.~1--40, 2009.
		
		\bibitem{Shor1994}
		P.\,W.\ Shor,
		``Algorithms for quantum computation: Discrete logarithms and factoring,''
		in \textit{Proc.\ 35th Annual Symposium on Foundations of Computer Science},
		S.~124--134, 1994.
		
		\bibitem{Bos2018}
		J.\ Bos u.\,a.,
		``CRYSTALS -- Kyber: A CCA-secure module-lattice-based KEM,''
		in \textit{Proc.\ IEEE EuroS\&P}, S.~353--367, 2018.
		
		\bibitem{Ducas2018}
		L.\ Ducas u.\,a.,
		``CRYSTALS -- Dilithium: A lattice-based digital signature scheme,''
		\textit{IACR Trans.\ Cryptographic Hardware and Embedded Systems},
		Bd.~2018, Nr.~1, S.~238--268, 2018.
		
		\bibitem{Cook1971}
		S.\,A.\ Cook,
		``The complexity of theorem-proving procedures,''
		in \textit{Proc.\ 3rd Annual ACM STOC}, S.~151--158, 1971.
		
		\bibitem{Khalil2002}
		H.\,K.\ Khalil,
		\textit{Nonlinear Systems}, 3.~Aufl.\ Upper Saddle River, NJ:
		Prentice Hall, 2002.
		
		\bibitem{Nipkow2002}
		T.\ Nipkow, L.\,C.\ Paulson und M.\ Wenzel,
		\textit{Isabelle/HOL: A Proof Assistant for Higher-Order Logic},
		Berlin: Springer, 2002.
		
		\bibitem{Lamport2002}
		L.\ Lamport,
		\textit{Specifying Systems: The TLA$^+$ Language and Tools},
		Boston: Addison-Wesley, 2002.
		
		\bibitem{Epp2026edfplus}
		S.\ Epp,
		\textit{Optimaler Scheduling-Algorithmus EDF$^+$: Prioritätsstrafe und
		stochastische Analyse bei 40 Tasks}, 2026. Verfügbar auf
		\texttt{github.com/hjstephan86/science/tree/main/science/edfplus}
		
		\bibitem{Epp2026pymca8}
		S.\ Epp,
		\textit{PYMCA-8: Python Multi-Core Architecture 8}, 2026. Verfügbar auf
		\texttt{github.com/hjstephan86/science/tree/main/science/pymca8}
		
		\bibitem{Epp2026msubgraph}
		S.\ Epp,
		\textit{The Hierarchical Subgraph-DFS Algorithm}, 2026. Verfügbar auf
		\texttt{github.com/hjstephan86/science}
		
		\bibitem{Epp2026subgraph}
		S.\ Epp,
		\textit{Der Subgraph Algorithmus: Graph-Profile und $O(n^3)$-Isomorphismus}, 2026. Verfügbar auf \texttt{github.com/hjstephan86/science}
		
		\bibitem{Epp2026boolmm}
		S.\ Epp,
		\textit{Boolesche Matrixmultiplikation in $O(n^{2.37})$}, 2026. Verfügbar auf \texttt{github.com/hjstephan86/science/tree/main/science/bool-mm}
		
		\bibitem{Epp2026archv}
		S.\ Epp,
		\textit{Das stochastische Archiv-Problem: Optimaler Batch-Timer}, 2026. Verfügbar auf
		\texttt{github.com/hjstephan86/science/tree/main/science/archv}
		
		\bibitem{Epp2026sysstate}
		S.\ Epp,
		\textit{Dynamische Systemzustandsklassifikation KS-I bis KS-IV}, 2026. Verfügbar auf
		\texttt{github.com/hjstephan86/science/tree/main/science/sysstate}
		
		\bibitem{Epp2026dusgraph}
		S.\ Epp,
		\textit{Dateisystem-Deduplizierung via Subgraph-Isomorphismus}, 2026. Verfügbar auf
		\texttt{github.com/hjstephan86/science/tree/main/science/dusgraph}
		
		\bibitem{Epp2026autsre}
		S.\ Epp,
		\textit{AUTOSAR OTA-Update-Architektur für Software-Defined Vehicles}, 2026. Verfügbar auf
		\texttt{github.com/hjstephan86/science/tree/main/science/autsre}
		
		\bibitem{Epp2026autsrepy}
		S.\ Epp,
		\textit{AUTOSAR Python-Treiber-Framework}, 2026. Verfügbar auf
		\texttt{github.com/hjstephan86/science/tree/main/science/autsrepy}
		
		\bibitem{Epp2026ethercate}
		S.\ Epp,
		\textit{EHAE: EtherCAT Hardware Abstraction Extension}, 2026. Verfügbar auf
		\texttt{github.com/hjstephan86/science/tree/main/science/ethercate}
		
		\bibitem{Epp2026liionp}
		S.\ Epp,
		\textit{KI-gestütztes Power Management Model für Lithium-Ionen-Akkumulatoren}, 2026. Verfügbar auf \texttt{github.com/hjstephan86/science/tree/main/science/liionp}
		
		\bibitem{Epp2026iono}
		S.\ Epp,
		\textit{Ionische Feldeffekttransistoren (IoFET) und Sigmoid-DVFS}, 2026. Verfügbar auf
		\texttt{github.com/hjstephan86/science/tree/main/science/iono}
		
		\bibitem{Epp2026pqc}
		S.\ Epp,
		\textit{Post-Quanten-Kryptographie: LWE und Gitter-basierte Adressierung}, 2026. Verfügbar auf
		\texttt{github.com/hjstephan86/science/tree/main/science/pqc}
		
		\bibitem{Epp2026physd}
		S.\ Epp,
		\textit{Physikalische Grundlagen digitaler Schaltkreise}, 2026. Verfügbar auf
		\texttt{github.com/hjstephan86/science/tree/main/science/physd}
		
		\bibitem{Epp2026signalth}
		S.\ Epp,
		\textit{Signaltheorie: Fourier-Transformation und Abtasttheorem}, 2026. Verfügbar auf
		\texttt{github.com/hjstephan86/science/tree/main/science/signalth}
		
		\bibitem{Epp2026mobde}
		S.\ Epp,
		\textit{Mobile Betriebssystem-Dateisystem-Architektur}, 2026. Verfügbar auf
		\texttt{github.com/hjstephan86/science/tree/main/science/mobde}
		
		\bibitem{Sha1990}
		L.\ Sha, R.\ Rajkumar und J.\,P.\ Lehoczky,
		``Priority inheritance protocols: An approach to real-time synchronization,''
		\textit{IEEE Trans.\ Computers}, Bd.~39, Nr.~9, S.~1175--1185, 1990.
		
		\bibitem{Baruah1990}
		S.\,K.\ Baruah, L.\,E.\ Rosier und R.\,R.\ Howell,
		``Algorithms and complexity concerning the preemptive scheduling of
		periodic, real-time tasks on one processor,''
		\textit{Real-Time Systems}, Bd.~2, Nr.~4, S.~301--324, 1990.
		
		\bibitem{Vestal2007}
		S.\ Vestal,
		``Preemptive scheduling of multi-criticality systems,''
		in \textit{Proc.\ 28th IEEE RTSS}, S.~239--243, 2007.
		
		\bibitem{Burns2013}
		A.\ Burns und R.\ Davis,
		\textit{Mixed Criticality Systems -- A Review},
		v5.0, Univ.\ of York, 2013.
		
		\bibitem{Feller1968}
		W.\ Feller,
		\textit{An Introduction to Probability Theory and Its Applications},
		Bd.~1, 3.~Aufl.\ New York: Wiley, 1968.

		\bibitem{Klein2009}
		G.\ Klein, K.\ Elphinstone, G.\ Heiser u.\,a.,
		``seL4: Formal verification of an OS kernel,''
		in \textit{Proc.\ 22nd ACM Symposium on Operating Systems Principles (SOSP)},
		S.~207--220, 2009.

		\bibitem{Epp2026vadis}
		S.\ Epp,
		\textit{Skalierbare Vektorisierung und KI-gestützte Strukturierung heterogener
		Großdaten: Ein Framework für On-Demand-Darstellung (VADIS)}, 2026. Verfügbar auf
		\texttt{github.com/hjstephan86/science/tree/main/science/vadis}
		
	\end{thebibliography}
	
\end{document}
